\documentclass[11pt,letterpaper,reqno]{amsart}
\usepackage[T1]{fontenc}
\usepackage{lmodern}
\usepackage[margin=1in]{geometry}
\usepackage{amsmath,amssymb,amsthm,mathtools,mathrsfs}
\usepackage{booktabs,microtype}
\usepackage[colorlinks=true,linkcolor=red,citecolor=red,urlcolor=blue,
  pdfauthor={ChatGPT},pdftitle={Holomorphic twists and an ordinary Fried obstruction},
  pdfsubject={A revised exposition of a proof draft},
  pdfkeywords={Fried conjecture, Ruelle zeta function, dynamical torsion, hyperbolic manifold}]{hyperref}

\newtheorem{theorem}{Theorem}[section]
\newtheorem{lemma}[theorem]{Lemma}
\newtheorem{proposition}[theorem]{Proposition}
\newtheorem{corollary}[theorem]{Corollary}
\theoremstyle{remark}
\newtheorem{remark}[theorem]{Remark}
\newcommand{\LL}{\mathcal L}
\newcommand{\Id}{\operatorname{Id}}
\newcommand{\Res}{\operatorname{Res}}
\newcommand{\tr}{\operatorname{tr}}
\newcommand{\str}{\operatorname{str}}
\newcommand{\Hol}{\operatorname{Hol}}
\newcommand{\Det}{\operatorname{Det}}
\newcommand{\ran}{\operatorname{ran}}
\newcommand{\vol}{\operatorname{vol}}
\newcommand{\cH}{\mathcal H}
\newcommand{\cU}{\mathcal U}
\newcommand{\cC}{\mathcal C}

\title[Holomorphic twists and a Fried obstruction]
{Holomorphic twists and an ordinary Fried obstruction\newline near hyperbolic metrics}
\author{ChatGPT}
\date{11 October 2026; research proof draft.}
\keywords{Fried conjecture, Ruelle zeta function, dynamical torsion, contact Anosov flow, holomorphic character deformation}

\begin{document}
\begin{abstract}
We construct an obstruction to the ordinary Fried identity arbitrarily
close to a hyperbolic three-metric. Along negatively curved metrics
converging to it and acyclic unitary characters tending to the trivial
character, the Ruelle zeta value divided by analytic torsion squared
tends to zero. A resonance crossing at quadratic character scale drives
the limit, controlled by uniform character holomorphy and finite-cut
torsion variation.
\end{abstract}
\maketitle

\noindent\textbf{AI statement.}
Simon Becker reports that the original argument was obtained from a
single prompt containing only the conjecture, submitted to
GPT-6.1 Sol with Ultra reasoning in Codex. Subsequent prompts revised
the exposition and references.

\section{Introduction}

Let $(\Sigma,g)$ be a closed oriented negatively curved Riemannian manifold of dimension
$n>2$, and let $\chi$ be an acyclic unitary representation of
$\pi_1(\Sigma)$. Write $\Id$ for the identity endomorphism and
$\mathcal G_g^{\#}$ for the primitive oriented periodic orbits of the
geodesic flow; both orientations of each closed geodesic are included.
For $\gamma\in\mathcal G_g^{\#}$, let $\ell_g(\gamma)$ be its length
and $\chi(\gamma)$ its holonomy. For a complex variable $s$, the Ruelle
zeta function in the convention used here is
\begin{equation}\label{Euler}
 Z_{g,\chi}(s)=\prod_{\gamma\in\mathcal G_g^{\#}}
 \det\bigl(\Id-\chi(\gamma)e^{-s\ell_g(\gamma)}\bigr),
 \qquad \operatorname{Re}s\gg1.
\end{equation}
Let $T_\chi(\Sigma)$ denote the Ray--Singer analytic torsion of the
local system defined by $\chi$. Conjecture~9 in Zworski's
survey~\cite{Zw} asserts
\begin{equation}\label{Friedidentity}
 |Z_{g,\chi}(0)|^{(-1)^{n-1}}=T_\chi(\Sigma)^2.
\end{equation}
In dimension three,
the exponent in \eqref{Friedidentity} is $+1$. The spectral coordinate
used in the survey is a complex variable $\lambda$; its zeta function,
denoted by $\zeta_\chi$, is related to ours by
\[
 \zeta_\chi(\lambda)=Z_{g,\chi}(-i\lambda).
\]

The value at zero requires separate attention when a resonance reaches
the origin. A zero resonance can contribute both numerator and
denominator factors to the dynamical determinant. The question is
therefore the value of the specialized meromorphic function, after its
spectral factors have been cancelled.


We write $\mathbf1$ for the trivial character. The conclusion of the
argument is the following.
\begin{theorem}\label{mainresult}
There exist a closed oriented hyperbolic three-manifold $(\Sigma,g_0)$,
negatively curved metrics $g_j\to g_0$ in $C^\infty$, and acyclic unitary
line characters $\chi_j\to\mathbf1$ such that $Z_{g_j,\chi_j}$ is
holomorphic at zero and
\begin{equation}\label{mainratio}
 0\leq\frac{|Z_{g_j,\chi_j}(0)|}{T_{\chi_j}(\Sigma)^2}
 \longrightarrow0.
\end{equation}
Consequently the asserted identity \eqref{Friedidentity} fails for every
sufficiently large $j$.
\end{theorem}

The decisive steps are uniform holomorphy on a fixed character disk and the extension
of the finite-cut variation calculation to a cut containing nonzero
resonances. The latter extension is proved explicitly in
Section~\ref{sec:normalization}.

The historical results fix both the benchmark and the scope of the
question. Ray and Singer~\cite{RS} defined analytic torsion from the
twisted Hodge spectrum and conjectured its equality with
Reidemeister--Franz torsion. Let $\Delta_{\chi,k}$ be the positive
twisted Hodge Laplacian on degree-$k$ forms. Write
$\zeta_{\Delta_{\chi,k}}$ for its spectral zeta function, formed from
the positive eigenvalues and continued to zero, and define
$\det{}'(\Delta_{\chi,k})=\exp(-\zeta_{\Delta_{\chi,k}}'(0))$;
the prime excludes zero eigenvalues. Our convention is
\begin{equation}\label{lit:RS}
 \log T_\chi(\Sigma)
  =\frac12\sum_{k=0}^{\dim\Sigma}(-1)^k k\,
       \zeta_{\Delta_{\chi,k}}'(0),
 \qquad
 T_\chi(\Sigma)=\prod_k
       \det{}'(\Delta_{\chi,k})^{(-1)^{k+1}k/2}.
\end{equation}
For an acyclic unitary system this is independent of the Riemannian
metric. Cheeger~\cite{Cheeger} and M\"uller~\cite{Muller} established
the equality with combinatorial torsion; M\"uller later extended the
comparison to unimodular representations~\cite{Muller93}. These results
make torsion a fixed benchmark along the metric deformation. They do
not identify a dynamical zeta value.

Fried's hyperbolic theorem~\cite{Fried} supplies that identification for
closed oriented hyperbolic manifolds with acyclic orthogonal
coefficients. For an odd-dimensional base, the ordinary oriented-orbit
product is regular and nonzero at the origin, with modulus equal to
the square of the base torsion. His nonacyclic theorem instead
determines a leading Laurent term. A complex unitary local system can
be treated by realification: its real torsion is the square of its
complex torsion, and its real zeta function is the product of the
complex zeta functions for the representation and its conjugate.
Reversal of oriented geodesics identifies these latter functions.
Taking the positive square root gives the complex-unitary
normalization used here. This translation supplies the hyperbolic
calibration, without any uniform statement near the trivial character.

The square also reflects the distinction between the base $\Sigma$ and
the flow manifold $M=S\Sigma$, its unit tangent sphere bundle. Let
$\pi:M\to\Sigma$ be the bundle projection, and use
$H^\bullet(\Sigma;\chi)$ for cohomology with coefficients in the local
system defined by $\chi$. For an oriented base of dimension at least three, the
sphere-bundle homotopy sequence identifies the fundamental groups of
$\Sigma$ and $S\Sigma$. Pullback preserves acyclicity, and in odd base
dimension the torsion comparison is
\begin{equation}\label{lit:sphere-torsion}
 H^\bullet(\Sigma;\chi)=0
   \quad\Longleftrightarrow\quad
 H^\bullet(S\Sigma;\pi^*\chi)=0,
 \qquad
 T_{\pi^*\chi}(S\Sigma)=T_\chi(\Sigma)^2.
\end{equation}
See \cite[Propositions 5.1--5.2]{ShenSurvey}. General flow formulations
therefore compare with torsion on $S\Sigma$, whereas
\eqref{Friedidentity} is stated on the base. Here these manifolds have
dimensions five and three, respectively.

Fried's later formulations have related but different hypotheses.
\emph{Lefschetz formulas for flows}~\cite{Fried87} develops the locally
homogeneous setting with orbit indices; the conjecture on page~66 is
stated in that setting. His analytic-flow paper~\cite{Fried95} proposes,
on page~181, a torsion identity for negatively curved real analytic
geodesic flows using matrix coefficients on the unit tangent bundle
and orientation factors. Translating coefficient assumptions, orbit
weights, and torsion conventions is necessary before identifying
these statements. The precise smooth, unitary, ordinary-value assertion
addressed by Theorem~\ref{mainresult} is Zworski's Conjecture~9.

The locally symmetric literature gives a larger positive class.
Moscovici--Stanton~\cite{MoSt} developed a trace-formula proof. Let $G$
be a connected linear real reductive group, $K$ a maximal compact
subgroup, and $\Gamma\subset G$ a torsion-free cocompact discrete
subgroup. Shen~\cite{Shen} established the unitary theorem for closed
odd-dimensional locally symmetric spaces $\Gamma\backslash G/K$.
His leading Laurent expansion has order zero and
torsion-squared coefficient when the bundle is acyclic. The proof uses
semisimple orbital integrals and representation theory, and completes
an argument in the earlier work. It covers the hyperbolic starting
metric, while the conformal family below leaves the locally symmetric
class. Shen's zeta and torsion conventions are both reciprocal to
those used here~\cite[(1-4)--(1-5)]{Shen}; both must be translated.
Locally symmetric formulas also retain Fuller-index weights when
closed geodesics occur in families.

Bunke--Olbrich~\cite{BObook,BO} connect Selberg functions to elliptic
operators on compact rank-one locally symmetric spaces. In the
odd-dimensional hyperbolic case, their determinant formulas express
Selberg factors through regularized elliptic determinants and explicit
nonvanishing exponential factors. The middle-degree Weyl pair requires
a graded determinant. These formulas identify local zeros and their
multiplicities in the appendix; Fried's theorem separately fixes the
modulus on the hyperbolic axis. In particular, local determinant
enumeration and global torsion normalization play different roles.

The continuation theory explains why the remaining issue is a special
value. Rugh~\cite{Rugh} used generalized Fredholm determinants for
analytic hyperbolic dynamics in dimension three, and
Fried~\cite{Fried95} treated analytic hyperbolic flows in arbitrary
dimension. Giulietti--Liverani--Pollicott~\cite{GLP} obtained
meromorphic continuation to the whole plane for smooth Anosov flows
using transfer operators on anisotropic spaces of currents. Their
erratum~\cite{GLPErr} corrects a contact spectral-gap argument without
affecting meromorphic continuation. None of these continuation results
alone ensures regularity, nonvanishing, or a torsion identity at zero.

The microlocal construction is especially suited to the present
deformation. Faure--Sj\"ostrand~\cite{FS} built anisotropic Sobolev
spaces with discrete generator spectrum in suitable half-planes;
Dyatlov--Zworski~\cite{DZ16} combined Fredholm theory, propagation,
and the flat trace to obtain a microlocal continuation proof.
Generalized resonant states can be distributions, and the ordinary
divisor is an alternating combination of their multiplicities.
Guedes Bonthonneau~\cite{Bon} constructed common anisotropic spaces
for nearby Anosov fields, giving continuous finite clusters and smooth
individual resonances away from collisions. At a collision the natural
objects are the contour projector and characteristic polynomial.
Cluster continuity alone neither selects smooth roots nor controls a
coefficient relative to a denominator tending to zero.

Dang--Rivi\`ere~\cite[Theorem 2.1]{DR} give the relevant topological
interpretation. For a flat vector bundle $E$ with connection $\nabla$
over $M$, let $d_\nabla$ be its flat de Rham differential and
$C^\bullet(0)$ the full generalized zero-resonance complex.
Projection onto that complex induces
\begin{equation}\label{lit:resonant-cohomology}
 H^k\bigl(C^\bullet(0),d_\nabla\bigr)
       \simeq H^k(M;E).
\end{equation}
This is a statement about the full de Rham complex, including
generalized states. It does not identify each horizontal resonance
space with a cohomology group. In particular,
\[
 H^\bullet(M;E)=0
 \quad\Longrightarrow\quad C^\bullet(0)\text{ is exact},
 \qquad
 C^\bullet(0)\text{ exact}\ \not\Longrightarrow\ C^\bullet(0)=0.
\]
A nonzero exact complex can contribute torsion when a resonance reaches
zero. This distinction is used in the crossing argument rather than
assuming horizontal exactness or semisimplicity.

Low-dimensional results give useful precedents with different ambient
dimensions. Write $\chi_{\rm top}(\Sigma)$ for the Euler characteristic.
Dyatlov--Zworski~\cite{DZ17} proved that the untwisted
surface-geodesic zeta function has order $-\chi_{\rm top}(\Sigma)$ at zero, and
also treated contact Anosov flows on three-dimensional flow manifolds.
Ceki\'c--Paternain~\cite{CP} studied volume-preserving flows in that
dimension, including unitary twists, twisted semisimplicity, and an
untwisted time-change example with a zero Jordan block. Their later
work~\cite{CP25} treats dissipative flows and Gaussian thermostats with
changed untwisted zeta order. An arbitrary time change or thermostat
need not be a Riemannian geodesic flow. These examples illuminate the
role of generalized multiplicities, but do not themselves supply an
acyclic unitary obstruction on the five-dimensional flow manifold here.

Dang--Guillarmou--Rivi\`ere--Shen~\cite{DGRS} prove local constancy of
the ordinary nonzero value where zero is absent from the full resonant
spectrum. Their Theorem~3 specializes this to geodesic flow over a
hyperbolic three-manifold: for a fixed acyclic unitary representation,
Selberg factorization excludes zero, and Fried's identity persists for
nearby Anosov fields. The representation is fixed before the
neighborhood is chosen. Denote the resulting metric neighborhood of
$g_0$ by $\mathcal V_\chi$. In schematic quantifier form, the conclusion is
\begin{equation}\label{lit:quantifiers}
 \text{for each fixed acyclic unitary }\chi,
 \quad \exists\,\mathcal V_\chi\ni g_0:
 \quad |Z_{g,\chi}(0)|=T_\chi(\Sigma)^2
       \quad(g\in\mathcal V_\chi).
\end{equation}
It does not assert one neighborhood valid for all characters tending
to the nonacyclic trivial character. The simultaneous limit in
Theorem~\ref{mainresult} therefore lies outside that uniform conclusion.

Ceki\'c--Delarue--Dyatlov--Paternain (CDDP)~\cite{CDDP} establish the
untwisted instability that supplies the geometric input. On a
hyperbolic three-manifold with first Betti number
$b_1=\dim H^1(\Sigma;\mathbb R)$, the horizontal
generalized zero ranks are $(1,2b_1,2b_1+2,2b_1,1)$. Under their
nondegeneracy condition, small nonzero perturbations have semisimple
zero ranks $(1,b_1,b_1+2,b_1,1)$; their conformal argument makes this
condition generic. Accordingly the ordinary vanishing order changes
from $4-2b_1$ to $4-b_1$. For $b_1=1$, the initial ranks are precisely
$(1,2,4,2,1)$. These are nonacyclic untwisted statements. The present
construction must additionally follow the twisted cluster and control
its specialized value at a character-dependent crossing.

Chaubet--Dang~\cite{CD} account for zero resonances by defining
dynamical torsion, denoted by $\tau_{\rm dyn}$. Equation numbers for
this reference follow the linked author version~v2. Write $Z$ for the
chosen Ruelle function, $m$ for its order at zero, and $a\ne0$ for
its leading Laurent coefficient,
so that $Z(s)=a\,s^m(1+O(s))$. For a contact form $\alpha$ on a
manifold of odd dimension $2q+1$, so that $q=(\dim M-1)/2$, let
$\Gamma_\alpha$ denote the induced
chirality, a degree-reversing involution; its explicit formula is given
in Section~\ref{sec:normalization}. In the acyclic case, write
$\tau(C^\bullet(0),\Gamma_\alpha)$ for the corresponding refined torsion
of the finite zero complex, and put
\[
 Q_0=\sum_{k=0}^{q}(-1)^k(q+1-k)\dim C^k(0).
\]
Their zero-cut comparison reads
\begin{equation}\label{lit:zero-correction}
 \tau_{\rm dyn}
   =(-1)^{Q_0}a^{(-1)^q}
        \tau\bigl(C^\bullet(0),\Gamma_\alpha\bigr).
\end{equation}
See \cite[(6.7)]{CD}; here $q=2$, hence the exponent of $a$ is $+1$.
Their construction is independent of the cut, invariant under contact
Anosov deformations, and holomorphic on the acyclic representation
locus; its quotient by suitable Turaev torsion is locally constant.
The ordinary value $Z(0)$, the leading coefficient $a$, and the
corrected torsion are distinct objects. Even when $m=0$, a nontrivial
zero complex can prevent replacing \eqref{lit:zero-correction} by
$\tau_{\rm dyn}=Z(0)$.

Bismut--Shen~\cite{BS} extend this perspective to a canonical
nonzero Fried section of the cohomological determinant line,
denoted by $\det H^\bullet(M;E)$, for general Anosov
fields, without excluding zero resonances. Their section is $C^1$ in
families and agrees with the contact construction. Gauss--Manin
flatness requires the bundle to be flat on the total fibration; a
family of fiberwise flat connections alone does not give that
hypothesis. Their norm-one identity for unimodular coefficients remains a modified conjecture.
Thus extending a regularized torsion through a finite spectral cut
already has substantial precedents. The scalar ordinary-value
calculation at the exceptional crossing is the issue here.

Recent generic results sharpen this distinction. In its March 2026
revision, Humbert--Tao~\cite{HT26} prove a fixed-representation
dichotomy on a connected open set of Anosov metrics: semisimplicity and
closedness of all horizontal zero states either hold on an open dense
set or hold nowhere. In the first case the dimensions are determined
by twisted cohomology. Their Corollary~1.2 extends the untwisted
generic description to the Anosov component containing a hyperbolic
three-metric. Their surface results also treat generic nonunitary
representations and produce exceptional acyclic twists with zero
Jordan blocks. Generic conclusions for fixed twists permit exceptional
metric loci; they do not give uniform control as a character approaches
the trivial representation, or assert the ordinary value at every
exceptional metric.

Other coefficient and flow classes require additional translations.
For nonunitary twists on odd-dimensional hyperbolic manifolds,
Spilioti~\cite{Spilioti25} compares the Ruelle value with
complex-valued torsions for acyclic representations sufficiently close
to the acyclic unitary locus. Shen~\cite{ShenComplex} develops the
locally symmetric comparison with Cappell--Miller torsion and, more
generally, leading terms. These conclusions concern complex torsion
and locally symmetric metrics; their neighborhoods do not include a
uniform passage to a nonacyclic character. J\'ez\'equel--Zung~\cite{JZ24}
prove a Fried identity for transitive smooth pseudo-Anosov flows on
closed three-dimensional flow manifolds, with acyclic unitary
coefficients and explicit monodromy conditions at singular orbits
(or a monodromy condition when there are no singular orbits).
Shen--Yu~\cite{ShenYu} relate Morse--Smale dynamics, Milnor metrics,
and torsion, with connecting trajectories contributing to the
comparison. These results broaden the dynamical setting but do not
cover the five-dimensional negatively curved geodesic flow considered
here.

The proposed obstruction consequently concerns a precise remaining
distinction: invariance of corrected dynamical torsion does not by
itself identify the ordinary specialized zeta value at a nonzero exact
zero complex. The argument below uses the published hyperbolic
calibration, perturbation framework, cohomological theorem, and
untwisted crossing mechanism. Its additional work is the uniform
complex character disk and the finite-cut normalization estimate at
quadratic character scale. Those estimates, rather than continuation
or regularization alone, are what must establish
Theorem~\ref{mainresult}.

For the construction, let $b_1(\Sigma)=\dim H^1(\Sigma;\mathbb R)$
be the first Betti number and choose the hyperbolic base so that
$b_1(\Sigma)=1$. Let $\omega$ be a nonzero real harmonic one-form,
$d$ the exterior differential, and $b$ a smooth real conformal
direction on $\Sigma$. For a small real metric parameter $\tau\geq0$
and a complex character parameter $z$, put
\[
 g_\tau=e^{-2\tau b}g_0,\qquad d_z=d+iz\omega\wedge,\qquad
 \chi_z(\gamma)=\exp\left(-iz\int_\gamma\omega\right),\qquad
 Z(s,z,\tau)=Z_{g_\tau,\chi_z}(s).
\]
We use $\theta$ for real values of $z$, so that $\chi_\theta$ is
unitary. Denote the two small horizontal degree-one roots of the
negative twisted Lie derivative by $c(z,\tau)$ and $t(z,\tau)$.
Set $\kappa=\|\omega\|_{L^2(\Sigma,g_0)}^2/
\operatorname{vol}_{g_0}(\Sigma)>0$, and denote by $r>0$ the initial
metric derivative of the companion root $t(0,\tau)$ for the chosen
direction. At the hyperbolic metric the two roots coincide, while the
untwisted metric variation moves the companion root:
\begin{equation}\label{mechanism}
 \begin{aligned}
 c(\theta,0)=t(\theta,0)
   &=-\tfrac12\kappa\theta^2+O(\theta^4),\qquad \kappa>0,\\
 t(0,\tau)&=r\tau+o(\tau),\qquad r>0.
 \end{aligned}
\end{equation}
Write $\sigma_\theta>0$ for a crossing parameter, characterized by
$t(\theta,\sigma_\theta)=0$; it satisfies
$\sigma_\theta\sim\kappa\theta^2/(2r)$. After the degree-one factors
are removed, the degree-two characteristic polynomial has a residual
quadratic factor $Q(s,z,\tau)$. Denote its linear coefficient by
$-e_1(z,\tau)$ and its constant coefficient by $e_2(z,\tau)$, so that
$Q(s,z,\tau)=s^2-e_1(z,\tau)s+e_2(z,\tau)$. Let $G(s,z,\tau)$ be the
holomorphic nonvanishing factor remaining after extraction of the
finite resonance divisor, constructed in Section~\ref{sec:normalization}.
At a crossing, exactness of the full zero complex forces $e_2=0$.
With $c$ and $e_1$ evaluated at $(\theta,\sigma_\theta)$, the ordinary
divisor becomes
\[
 Z(s,\theta,\sigma_\theta)
   =G(s,\theta,\sigma_\theta)(s-c)(s-e_1).
\]
The three estimates responsible for \eqref{mainratio} are
\begin{equation}\label{threeestimates}
 \frac{G(0,\theta,\sigma_\theta)}{G(0,\theta,0)}\longrightarrow1,
 \qquad
 \frac{c(\theta,\sigma_\theta)}{c(\theta,0)}\longrightarrow1,
 \qquad
 \frac{e_1(\theta,\sigma_\theta)}{c(\theta,0)}\longrightarrow0.
\end{equation}
Section~\ref{sec:spectral} constructs the common cut, proves the
coefficient estimates, and produces the crossing.
Section~\ref{sec:normalization} controls the fixed-cut normalization
and specializes the ordinary zeta function to prove
Theorem~\ref{mainresult}. The appendix records the geometric and
spectral inputs in their source conventions.

\section{The spectral cut and the crossing}\label{sec:spectral}
\subsection{Geometric setup and hyperbolic normalization}\label{sec:geometry}
Let $(\Sigma,g_0)$ be a closed connected oriented hyperbolic
three-manifold with $b_1=1$, and let $\omega\neq0$ be its real harmonic
one-form. Choose a smooth real function $b$ on $\Sigma$, to be specified
below, and a real metric parameter $\tau\geq0$ near zero. With $d$
denoting the exterior differential, set
\[
 d_z=d+iz\omega\wedge,\qquad g_\tau=e^{-2\tau b}g_0,
 \qquad z\in\mathbb C,
\]
and identify the unit bundles by radial scaling with $M=S_{g_0}\Sigma$:
\[
 F_\tau(x,v)=(x,e^{\tau b(x)}v),\qquad
 F_\tau:M\longrightarrow S_{g_\tau}\Sigma.
\]
Write $\pi:M\to\Sigma$ for the bundle projection. The character of
the flat connection on the trivial complex line bundle is
$\chi_z(\gamma)=\exp(-iz\int_\gamma\omega)$. On $M$ the same
notation $d_z$ denotes the pullback differential
$d+iz\pi^*\omega\wedge$.
Write $\alpha_\tau,X_\tau$ for the pulled-back geodesic contact form
and Reeb field. If $\iota_X$ denotes contraction with a vector field
$X$, the twisted Lie derivative and its negative are
\[
 L_{z,\tau}=\LL^{d_z}_{X_\tau}
       =d_z\iota_{X_\tau}+\iota_{X_\tau}d_z,
 \qquad B_{z,\tau}=-L_{z,\tau}.
\]
We suppress $(z,\tau)$ on operators when the parameter is clear.
A form is called horizontal when $\iota_{X_\tau}u=0$.
The spectral coordinate $s$ is an eigenvalue of $B$:
\[
 L+s=s\Id-B,\qquad (L+s)^{-1}=(s\Id-B)^{-1}.
\]
For comparison with CDDP's convention, $P=-iL=iB$ and $B=-iP$. For real $z=\theta\neq0$ sufficiently small, the
unitary character is acyclic. The construction of the base and the Hodge proof of acyclicity
are given in the appendix.


The conformal direction is chosen so that the untwisted horizontal
degree-one companion root has derivative $r>0$ at $\tau=0$.
Let $\Delta_{g_0}$ be the nonpositive scalar Laplacian and $Q_4$ the
self-adjoint periodized smoothing operator of CDDP's Theorem~5.
One explicit choice is $b=Q_4\Delta_{g_0}|\omega|^2$. The appendix
defines this operator in \eqref{app:Q4definition} and verifies the nonvanishing pairing
and the positive sign using CDDP's contact-orientation convention.
The actual primary input is CDDP (3.55)--(3.57), (4.22), (4.38),
Theorem 5, and Proposition A.1 \cite{CDDP}. With $b_1=1$, this is
the nondegeneracy condition in CDDP Theorem 3. Their compressed
$P=-iL$ variation is multiplied by $-i$ to obtain the variation of
$B=-iP$.


Fix $\lambda>0$ sufficiently small so that the positively oriented
circle $\gamma=\{|s|=\lambda\}$ contains no boundary resonances at
$(z,\tau)=(0,0)$ in any degree. The horizontal cut $\cU^k_{z,\tau}$ is the
sum of the generalized resonance spaces of $B_{z,\tau}$ in degree
$k$ inside $\gamma$, restricted to horizontal forms. Its ranks in
degrees $0,\ldots,4$ are
\begin{equation}\label{ranks}
 (1,2,4,2,1),
 \qquad\text{and its full ranks are }(1,3,6,6,3,1).
\end{equation}
The full ranks refer to degrees $0,\ldots,5$, with no horizontality
restriction. Define the compressed operators, scalar root, companion
root, and characteristic polynomials by
\[
 A_k(z,\tau)=B_{z,\tau}|_{\cU^k_{z,\tau}},\qquad
 c(z,\tau)=A_0(z,\tau),\qquad
 t(z,\tau)=\tr A_1(z,\tau)-c(z,\tau),
\]
\[
 P_k(s,z,\tau)=\det(s\Id-A_k(z,\tau)).
\]
The rank-one operator $A_0$ is identified with its scalar root $c$.
Let $\mu_0(\theta)$ be the small positive twisted scalar Hodge
eigenvalue, and put $c_0(\theta)=-1+\sqrt{1-\mu_0(\theta)}$,
using the square-root branch equal to one at the origin.
The Selberg enumeration in the appendix, including
its scalar denominator cancellation, gives
on the hyperbolic axis, for small real $\theta\neq0$,
\begin{equation}\label{hyper}
 c(\theta,0)=t(\theta,0)=c_0(\theta)
       =-1+\sqrt{1-\mu_0(\theta)}<0,
 \qquad
 \det(s-A_2(\theta,0))=(s-c_0)^3(s+c_0).
\end{equation}
These are the consequences of DGRS (7.6) and its Selberg
normalization used here \cite{DGRS}; they are not deduced from abstract
contact algebra.


For later reference,
\begin{equation}\label{quadratic}
 \mu_0(\theta)=\kappa\theta^2+O(\theta^4),\qquad
 c_0(\theta)=-\tfrac12\kappa\theta^2+O(\theta^4),
 \qquad
 \kappa=\frac{\|\omega\|_{L^2(\Sigma)}^2}
                  {\operatorname{vol}(\Sigma)}>0.
\end{equation}
Let $d^*$ denote the $g_0$-adjoint of $d$. The first derivative of
the scalar Laplacian annihilates the normalized constant, since
$d^*\omega=0$. The first correction to that
simple normalized eigenfunction is therefore zero. The coefficient
of $\theta^2$ in its eigenvalue is the expectation of $|\omega|^2$.
Conjugation under $\theta\mapsto-\theta$ makes this eigenvalue even.
Expanding the square root gives the second assertion.

\subsection{Uniform holomorphy and reversal symmetry}\label{sec:holom}
\begin{lemma}\label{holom}
Write $\Hol(D)$ for the holomorphic functions on a complex domain
$D$. After shrinking the parameter neighborhood, there are
$\varepsilon>0$ and one fixed $\rho>0$ such that the finite matrices
for the cut and their characteristic coefficients belong to
\[
 C^0\bigl([0,\varepsilon]_\tau;
             \Hol(\{|z|<\rho\})\bigr).
\]
The topology is locally uniform convergence
in the complex $z$ variable. Their ranks are \eqref{ranks}.
\end{lemma}
\begin{proof}
In each form degree, fix a finite collection of anisotropic spaces
\(\mathcal H_0,\ldots,\mathcal H_4\) for the untwisted metric family,
as in \cite[Appendix B.4--B.6]{CD}. Write
$R_{z,\tau}(s)=(L_{z,\tau}+s)^{-1}$ for the resolvent, suppressing
the form-degree index. For each degree and each of these levels, the
untwisted contour resolvent satisfies
\begin{equation}\label{xp:resolvent-bounds}
 \sup_{\substack{s\in\gamma\\0\leq\tau\leq\varepsilon}}
 \|R_{0,\tau}(s)\|_{\mathcal H_j\to\mathcal H_j}<\infty,
 \qquad
 R_{0,\tau}(s)=(L_{0,\tau}+s)^{-1}.
\end{equation}
The twist is an order-zero multiplication operator, as in
\cite[(8.9)]{CD}:
\begin{equation}\label{xp:twist-multiplier}
 L_{z,\tau}=L_{0,\tau}+zV_\tau,
 \qquad V_\tau=i\,\pi^*\omega(X_\tau),
 \qquad
 \sup_{\tau}\|V_\tau\|_{\mathcal H_j\to\mathcal H_j}<\infty.
\end{equation}
Choose one radius \(\rho>0\) so small that, simultaneously on this
finite list of spaces,
\begin{equation}\label{xp:uniform-neumann-smallness}
 \sup_{\substack{|z|<\rho,\ s\in\gamma\\0\leq\tau\leq\varepsilon}}
 \|zV_\tau R_{0,\tau}(s)\|_{\mathcal H_j\to\mathcal H_j}
 \leq\tfrac12.
\end{equation}
The factorization
\[
 L_{z,\tau}+s
   =(\Id+zV_\tau R_{0,\tau}(s))(L_{0,\tau}+s)
\]
then yields the convergent Neumann series
\begin{equation}\label{Neumann}
 R_{z,\tau}(s)
 =R_{0,\tau}(s)(\Id+zV_\tau R_{0,\tau}(s))^{-1}
 =\sum_{n=0}^{\infty}(-z)^n
       R_{0,\tau}(s)(V_\tau R_{0,\tau}(s))^n.
\end{equation}
This proves holomorphy in the same disk for every nearby metric and
gives uniform bounds for the twisted resolvent.

Metric continuity is used with a loss of one level. Let $C_j$ bound
the contour resolvent norms on both levels $j$ and $j+1$, uniformly
in the chosen parameter neighborhood. Then
\begin{align}\label{xp:resolvent-continuity}
 R_{z,\tau}(s)-R_{z,\tau'}(s)
 &=R_{z,\tau}(s)
       (L_{z,\tau'}-L_{z,\tau})R_{z,\tau'}(s),\\
 \|R_{z,\tau}(s)-R_{z,\tau'}(s)\|_{\mathcal H_{j+1}\to\mathcal H_j}
 &\leq C_j^2
       \|L_{z,\tau'}-L_{z,\tau}\|_{\mathcal H_{j+1}\to\mathcal H_j}
       \longrightarrow0.\notag
\end{align}
The last convergence is uniform for \(s\in\gamma\) and for \(z\)
in any smaller closed disk. Define the full-cut spectral projection
$\Pi_{z,\tau}$ and its first spectral moment $\mathsf T_{z,\tau}$
by contour integration:
\begin{equation}\label{xp:projector-moment}
 \Pi_{z,\tau}=\frac1{2\pi i}\int_\gamma R_{z,\tau}(s)\,ds,
 \qquad
 \mathsf T_{z,\tau}=\frac1{2\pi i}
                   \int_\gamma sR_{z,\tau}(s)\,ds
                 =B_{z,\tau}\Pi_{z,\tau}.
\end{equation}
They have the same holomorphic and continuous dependence.

The untwisted ranks persist by \cite[Appendix B.6]{CD}.  After
shrinking \(\rho\), the Neumann series also gives
\[
 \|\Pi_{z,\tau}-\Pi_{0,\tau}\|_{\mathcal H_j\to\mathcal H_j}<1,
\]
so the two finite-rank projections have equal ranks.  To pass from
operators to finite matrices, let $N$ be the cut rank in the chosen
form degree and choose fixed smooth sources
\(f_1,\ldots,f_N\) and fixed smooth dual tests
\(\ell_1,\ldots,\ell_N\), viewed as linear functionals on that
degree, such that
\begin{equation}\label{xp:frame-matrices}
 F_{ij}(z,\tau)=\ell_i(\Pi_{z,\tau}f_j),
 \qquad
 K_{ij}(z,\tau)=\ell_i(\mathsf T_{z,\tau}f_j),
 \qquad \det F(0,0)\neq0.
\end{equation}
Then \(\det F(z,\tau)\neq0\) nearby, the vectors
\(\Pi_{z,\tau}f_j\) form a frame, and the compressed operator has
matrix
\begin{equation}\label{xp:compressed-matrix}
 [B_{z,\tau}]_{\{\Pi f_j\}}=F(z,\tau)^{-1}K(z,\tau).
\end{equation}
This formula makes the asserted regularity of every characteristic
coefficient explicit.  For the horizontal degrees one first inserts
\(h_\tau=\Id-\alpha_\tau\wedge\iota_{X_\tau}\), which commutes
with the resolvent and identifies the horizontal bundles.  No
same-level operator-norm continuity in \(\tau\) is required.
\end{proof}


\begin{lemma}\label{even}
The characteristic coefficients in Lemma~\ref{holom} are even in
$z$ and real for real $z$. In particular $c$ is even and real on the
real locus.
\end{lemma}
\begin{proof}
The reversal argument applies to each horizontal degree separately.
For a primitive oriented orbit $\gamma$, let $\ell_\gamma$ be its
length, $\bar\gamma$ its reverse, $P_\gamma$ its full transverse
symplectic return map, and $I_\gamma=\int_{\pi\gamma}\omega$.
Write $\Lambda^k$ for the $k$th exterior power. In a right
half-plane, the degree-$k$ dynamical determinant $\mathscr D_k$ is
\begin{equation}\label{xp:degreewise-euler}
 \mathscr D_k(s,z,\tau)=\exp\!\left[
  -\sum_{\gamma\ {\rm primitive}}\sum_{j\geq1}
    \frac{e^{-sj\ell_\gamma}}{j}
    \frac{e^{-izjI_\gamma}\,
          \tr(\Lambda^kP_\gamma^j)}
         {|\det(\Id-P_\gamma^j)|}\right].
\end{equation}
Let $\varphi_t$ be the geodesic flow of $X_\tau$. The reversal
\(R(x,v)=(x,-v)\) satisfies, with $\sim$ denoting similarity of
linear maps,
\[
 R\varphi_t=\varphi_{-t}R,
 \qquad
 I_{\bar\gamma}=-I_\gamma,
 \qquad
 P_{\bar\gamma}\sim P_\gamma^{-1}.
\]
If $\mathbb J$ is the standard symplectic matrix in a symplectic
basis, then
\[
 (P_\gamma^j)^T\mathbb J P_\gamma^j=\mathbb J
 \quad\Longrightarrow\quad
 P_\gamma^{-j}=\mathbb J^{-1}(P_\gamma^j)^T\mathbb J.
\]
Consequently, for every \(k\),
\begin{equation}\label{xp:inverse-exterior-traces}
 \tr(\Lambda^kP_\gamma^{-j})
   =\tr(\Lambda^kP_\gamma^j),
 \qquad
 |\det(\Id-P_\gamma^{-j})|
   =|\det(\Id-P_\gamma^j)|.
\end{equation}
Pairing the two orientations therefore replaces the holonomy factor
by
\[
 e^{-izjI_\gamma}+e^{izjI_\gamma}
       =2\cos(zjI_\gamma).
\]
The paired Euler sum gives
\begin{equation}\label{xp:degreewise-symmetry}
 \mathscr D_k(s,-z,\tau)=\mathscr D_k(s,z,\tau),
 \qquad
 \mathscr D_k(\bar s,\theta,\tau)
       =\overline{\mathscr D_k(s,\theta,\tau)}
       \quad(\theta\in\mathbb R).
\end{equation}
These identities continue with the determinants.  Their zero
divisors are the degreewise resonance multisets, including algebraic
multiplicity.  The monic characteristic polynomial of the fixed,
conjugation-invariant cut is uniquely determined by that multiset;
it is therefore even in \(z\) and real for real \(z\).
\end{proof}

\subsection{The characteristic coefficients}
On the zero-twist axis $c(0,\tau)=0$. For real nonzero small $\theta$, the scalar
zero exclusion of DGRS Lemma 7.4 implies $c\neq0$. The horizontal
operators $D=d_z-\alpha_\tau\wedge L$ and $W=d\alpha_\tau\wedge$
satisfy
\begin{equation}\label{Cartan}
 [B,D]=0,\qquad D^2=WB.
\end{equation}
If $f$ is a nonzero scalar cut state with $Bf=cf$ and $c\ne0$, then
\[
 D^2f=cWf\ne0,\qquad B(Df)=cDf.
\]
Thus $Df\ne0$ and $c$ is also a degree-one eigenvalue. Therefore
\begin{equation}\label{P1}
 \det(s\Id-A_1)=(s-c)(s-t).
\end{equation}
The equality extends first to complex $z$ by holomorphy and then to
$z=0$ by continuity. From the primary first variation,
\begin{equation}\label{taxis}
 t(0,\tau)=r\tau+o(\tau),\qquad r>0.
\end{equation}
Only a derivative on this untwisted axis is used.


\begin{lemma}\label{scalarfactors}
On a smaller fixed disk there are functions $C$ and $T$, holomorphic in
$z$ and continuous in $\tau$ locally uniformly in $z$, such that
\begin{equation}\label{ctfactors}
 c(z,\tau)=z^2C(z,\tau),\qquad
 t(z,\tau)=t(0,\tau)+z^2T(z,\tau).
\end{equation}
Their values at the origin are
\begin{equation}\label{ctlimits}
 C(0,0)=T(0,0)=-\kappa/2.
\end{equation}
In particular $c(z,\tau)\ne0$ for sufficiently small complex $z\ne0$.
\end{lemma}
\begin{proof}
If $h$ is even and holomorphic, with $h(0,\tau)=0$, choose a radius
$0<a<\rho$. Its quotient by $z^2$ has the Cauchy representation
\begin{equation}\label{Cauchyquotient}
 q_h(z,\tau):=\frac{h(z,\tau)}{z^2}
 =\frac{1}{2\pi i}\int_{|w|=a}
       \frac{h(w,\tau)}{w^2(w-z)}\,dw,
 \qquad |z|<a,
\end{equation}
where the value at $z=0$ is removable. For a smaller radius $0<a'<a$,
\begin{equation}\label{quotientcontinuity}
 \sup_{|z|\leq a'}|q_h(z,\tau)-q_h(z,\tau')|
 \leq\frac{1}{a(a-a')}
   \sup_{|w|=a}|h(w,\tau)-h(w,\tau')|.
\end{equation}
Lemma~\ref{holom} makes the right side tend to zero as $\tau'\to\tau$.
Apply this to $h=c$ and $h=t-t(0,\tau)$. Evenness comes from
Lemma~\ref{even}, and both functions vanish at $z=0$. The hyperbolic
expansion \eqref{quadratic} gives \eqref{ctlimits}. Shrinking the
neighborhood once more, we may arrange
\[
 |C(z,\tau)+\kappa/2|<\kappa/4,
 \qquad |C(z,\tau)|\geq\kappa/4,
\]
which proves the scalar nonvanishing assertion.
\end{proof}


Whenever $A_1$ is invertible,
$D:\cU^1_{z,\tau}\to\cU^2_{z,\tau}$ is
injective, because
\[
 Du=0\ \Longrightarrow\ D^2u=WBu=0
 \ \Longrightarrow\ Bu=0\ \Longrightarrow\ u=0.
\]
Here the middle implication uses the contact Lefschetz isomorphism
$W:\cU^1_{z,\tau}\to\cU^3_{z,\tau}$. Since $[B,D]=0$, the image
of $D$ is a $B$-invariant subspace. Thus \eqref{P1} divides
$\det(s\Id-A_2)$ wherever $A_1$ is invertible. At every parameter,
let $Q$ be the monic quadratic quotient and $\mathcal R$ the
remainder in the polynomial division
\[
 P_2(s)=P_1(s)Q(s)+\mathcal R(s),\qquad
 \deg\mathcal R<2.
\]
The coefficients of $Q$ and $\mathcal R$ inherit holomorphy in $z$
and continuity in $\tau$.
For each fixed small $\tau>0$, $t(0,\tau)\neq0$, so its holomorphic
$z$-zeros are discrete unless it is identically zero, which it is not.
Lemma~\ref{scalarfactors} gives $c(z,\tau)\ne0$ for
$z\ne0$; at $\tau=0$ the hyperbolic identity \eqref{hyper} applies. Thus the invertible
locus is dense, and the remainder vanishes identically. Denoting the
two coefficients of $Q$ by $e_1(z,\tau)$ and $e_2(z,\tau)$, we obtain
\begin{equation}\label{P2}
 \det(s\Id-A_2)=(s-c)(s-t)Q(s),\qquad
 Q(s)=s^2-e_1s+e_2.
\end{equation}
Monic division preserves holomorphy in $z$ and continuity in $\tau$.
Write $c_0(z):=c(z,0)$ for the holomorphic continuation of the real
hyperbolic root. The identities in \eqref{hyper} extend from real
characters by holomorphy and give
\begin{equation}\label{qhyper}
 e_1(z,0)=0,\qquad e_2(z,0)=-c_0(z)^2.
\end{equation}
On the zero-twist axis CDDP Theorem 3 gives horizontal zero
multiplicities one in degree one and three in degree two, for all
small $\tau>0$. The additional middle root is the nonzero $t(0,\tau)$
carried by $D$. Consequently
\begin{equation}\label{qtrivial}
 e_1(0,\tau)=e_2(0,\tau)=0.
\end{equation}


\begin{lemma}[The relative estimate]\label{relative}
There is a function $A$, holomorphic in $z$ and continuous in $\tau$
locally uniformly on a fixed disk, such that
\begin{equation}\label{factors}
 c=z^2C,\qquad t=t(0,\tau)+z^2T,\qquad e_1=z^2A.
\end{equation}
The three quotient values satisfy
\begin{equation}\label{limits}
 C(0,0)=T(0,0)=-\kappa/2,\qquad A(0,0)=0.
\end{equation}
\end{lemma}
\begin{proof}
The first two factors have already been constructed in
Lemma~\ref{scalarfactors}, before the polynomial division. The
coefficient $e_1$ is now defined by \eqref{P2}; it is even and
holomorphic in $z$, and $e_1(0,\tau)=0$ by \eqref{qtrivial}.
Apply \eqref{Cauchyquotient}--\eqref{quotientcontinuity} with $h=e_1$.
The identity $e_1(z,0)=0$ implies $A(z,0)=0$, including its removable
value at the origin. This proves $A(0,0)=0$.
\end{proof}

\subsection{The crossing and the full zero complex}\label{sec:crossings}
Choose \(K>0\) such that \(|T|\leq K\).  On a sufficiently small
positive metric interval,
\[
 t(0,\tau)\geq\tfrac r2\tau,
 \qquad
 t(\theta,\tau)=t(0,\tau)+\theta^2T(\theta,\tau).
\]
For \(\tau_1=4K\theta^2/r\),
\[
 t(\theta,\tau_1)
       \geq\tfrac r2\tau_1-K\theta^2
       =K\theta^2>0,
 \qquad t(\theta,0)=c_0(\theta)<0.
\]
Continuity supplies a root
\begin{equation}\label{cross}
 t(\theta,\sigma_\theta)=0,
 \qquad 0<\sigma_\theta\leq4K\theta^2/r.
\end{equation}
More generally, every root in the chosen positive interval satisfies
\[
 \tfrac r2\sigma_\theta
     \leq\theta^2|T(\theta,\sigma_\theta)|
     \leq K\theta^2.
\]
Writing \(t(0,\tau)=\tau(r+\eta(\tau))\), where
\(\eta(\tau)\to0\), the root equation becomes
\begin{equation}\label{xp:crossing-asymptotic}
 \frac{\sigma_\theta}{\theta^2}
    =-\frac{T(\theta,\sigma_\theta)}{r+\eta(\sigma_\theta)}
    \longrightarrow\frac\kappa{2r}.
\end{equation}
No selection with uniqueness or monotonicity is needed.

At a crossing $c\ne0$, so the degree-one zero eigenvalue is simple.

For completeness, off the zero set of $t$ both scalar and horizontal
one-zero spaces vanish. DGRS Proposition 7.3, using pullback
acyclicity and contact splitting, excludes the full zero complex.
DGRS Theorem 2 then makes the ordinary zeta value constant on each
such connected interval. The proof below does not differentiate that
ordinary value or take its limit through a crossing.


Fix a crossing $(\theta,\sigma_\theta)$ and suppress its parameters.
Let $V^k$ be the horizontal generalized zero space and put
$m=\dim V^2$. Scalar exclusion and contact Lefschetz give
\[
 (\dim V^0,\ldots,\dim V^4)=(0,1,m,1,0).
\]
The degree-one zero is simple because its other cut eigenvalue is
\(c\neq0\).  In particular \(L|_{V^1}=0\); Lefschetz then gives
\(L|_{V^3}=0\) as well. Write $\alpha=\alpha_{\sigma_\theta}$
and use left wedging $jv=\alpha\wedge v$.
The full generalized zero complex has spaces
$C^k=V^k\oplus jV^{k-1}$, where $V^{-1}=V^5=0$.
Its dimensions and differential are
\begin{equation}\label{xp:full-zero-complex}
 (\dim C^0,\ldots,\dim C^5)=(0,1,m+1,m+1,1,0),
 \qquad
 d_\theta(u+jv)=Du+Wv+j(Lu-Dv).
\end{equation}
The resonant de Rham theorem makes this complex exact. Write $d_k$
for the restriction of $d_\theta$ to $C^k$ and $D_k$ for the restriction
of the horizontal differential $D$ to $V^k$. In degree one,
\(C^1=V^1\) and
\[
 d_1u=(Du,0),
 \qquad \ker d_1=0,
\]
so \(D_1:V^1\to V^2\) is injective.  In degree three,
\(W V^2\subset V^4=0\), \(DV^3\subset V^4=0\), and
\(L|_{V^3}=0\), giving
\[
 d_3(u+jv)=-jD_2v.
\]
Exactness in degree four, where \(C^4=jV^3\) and \(C^5=0\),
therefore makes \(D_2:V^2\to V^3\) surjective.  Finally,
\[
 D_2D_1=WB|_{V^1}=0,
 \qquad
 1=\dim\operatorname{im}D_1
    \leq\dim\ker D_2=m-1,
\]
so \(m\geq2\).  The nonzero root \(c\) occupies one dimension
of the rank-four middle cut, giving \(m\leq3\).  Since
\[
 P_2(s)=(s-c)sQ(s),
 \qquad Q(s)=s^2-e_1s+e_2,
\]
the order of \(s=0\) in \(P_2\) is exactly \(m\), and hence
\begin{equation}\label{two}
 e_2=0,
 \qquad
 m=3\Longleftrightarrow e_1=0,
 \qquad
 m=2\Longleftrightarrow e_1\neq0.
\end{equation}
This argument uses exactness of the full complex; it does not assume
exactness of the horizontal sequence.

\section{Fixed-cut torsion and the ordinary zeta value}\label{sec:normalization}
\subsection{The ordinary divisor}
Let $G(s,z,\tau)$ denote the complementary factor obtained by
removing the fixed spectral cut, in the convention of
\cite[(6.4)--(6.5)]{CD}. For each parameter it is holomorphic and
nonvanishing on the closed spectral disk. The full cut in degree $k$
is $\cU^k\oplus\alpha\wedge\cU^{k-1}$, so its determinant is
$P_kP_{k-1}$, with $P_{-1}=P_5=1$. The weighted full-cut determinant
telescopes:
\begin{align}\label{fulltelescoping}
 \mathfrak D_{\rm cut}(s)
 &=\prod_{k=0}^{5}(P_kP_{k-1})^{(-1)^k k}
   =\prod_{k=0}^{4}P_k^{(-1)^{k+1}}\notag\\
 &=\frac{P_1P_3}{P_0P_2P_4}
   =\frac{P_1^2}{P_0^2P_2}
   =\frac{s-t}{(s-c)Q(s)}.
\end{align}
The middle equality uses the contact Lefschetz identities $P_3=P_1$
and $P_4=P_0$. In flow dimension five the convention of
\cite[(6.5)]{CD} is $G=Z\mathfrak D_{\rm cut}$. Thus the ordinary
zeta divisor is
\begin{equation}\label{divisor}
 Z(s,z,\tau)=G(s,z,\tau)
   \frac{(s-c)(s^2-e_1s+e_2)}{s-t}.
\end{equation}
Only the following relative normalization is needed; no joint ordinary
zeta evaluation at $s=0$ is assumed.

\subsection{Variation on the full cut}
A dot denotes differentiation with respect to the metric parameter
$\tau$, and $\iota_Y$ denotes contraction by a vector field $Y$.
For a degree-preserving endomorphism $A$ of a graded finite complex
$C^\bullet$, its supertrace is
\[
 \str_{C^\bullet}A=\sum_k(-1)^k\tr_{C^k}A.
\]
\begin{lemma}[Arbitrary full cut variation]\label{cutvariation}
For fixed real nonzero small $\theta$, put
$C^\bullet_{\theta,\tau}=\operatorname{ran}\Pi_{\theta,\tau}$ for
the full cut \eqref{ranks}. Define
\[
 H(\theta,\tau)=
     \str_{C^\bullet_{\theta,\tau}}
       \bigl(\Pi_{\theta,\tau}\alpha_\tau\wedge
                                  \iota_{\dot X_\tau}\bigr).
\]
Then
\begin{equation}\label{Gvariation}
 \frac{G(0,\theta,\tau)}{G(0,\theta,0)}
       =\exp\left(\int_0^\tau H(\theta,v)\,dv\right).
\end{equation}
The finite trace $H$ has a bounded continuous extension to $\theta=0$.
Consequently the left side tends to one uniformly as $\tau\to0$
and $\theta\to0$ through real nonzero characters.
\end{lemma}
\begin{proof}
Hold $\theta$ fixed and abbreviate
$C^\bullet_\tau=C^\bullet_{\theta,\tau}$ and
$U^k_\tau=\cU^k_{\theta,\tau}$. When a metric subscript is omitted,
all operators are evaluated at the current parameter: in particular
$L=L_{\theta,\tau}$, $X=X_\tau$, $\alpha=\alpha_\tau$, and
$\Pi=\Pi_{\theta,\tau}$; also write $U^k=U^k_\tau$.
Write $W_\tau=d\alpha_\tau\wedge$.
The right-wedge contact splitting and the Lefschetz map are
\[
 C^k_\tau=U^{k-1}_\tau\wedge\alpha_\tau\oplus U^k_\tau,
 \qquad \ell_\tau u=u\wedge d\alpha_\tau,
 \qquad U^{-1}_\tau=U^5_\tau=0.
\]
We abbreviate $\ell=\ell_\tau$ at the current parameter.
For $f\in U^{k-1}$ and $g\in U^k$, the contact chirality
$\Gamma_\tau$ is the degree-reversing involution given, for
$k=0,1,2$, by the following formula, abbreviating $\Gamma=\Gamma_\tau$:
\begin{equation}\label{xp:contact-chirality}
 \Gamma(f\wedge\alpha+g)
       =\ell^{2-k}g\wedge\alpha+\ell^{3-k}f;
\end{equation}
in the higher degrees its action is the inverse. Here powers and
inverse powers of $\ell$ mean the appropriate contact Lefschetz
maps between horizontal degrees.

The cut complex is acyclic for every real nonzero small $\theta$.
Its zero part computes $H^\bullet(M;d_\theta)=0$. On a nonzero
resonance part, $\iota_XL^{-1}$ is a cochain contraction by Cartan's
identity. The fixed cut commutes with $d_\theta$, $\iota_{X_\tau}$,
wedging by $\alpha_\tau$, and $W_\tau$; the smooth inverse Lefschetz
maps commute as well. Thus the contact chirality $\Gamma_\tau$
preserves the cut. These statements hold on every generalized
resonance, not only on a zero block.

Chaubet--Dang initially choose a cut containing only zero at the
reference parameter. We give the trace calculation on the present
full cut. The smooth-family projector regularity of
\cite[Proposition 7.2 and Appendix B]{CD} supplies differentiable
finite cut matrices in \(\tau\), with a loss of anisotropic levels.
This fixed-character differentiability is used only for the torsion
variation formula.

For an acyclic finite complex $C^\bullet$ with chirality $\Gamma$,
$\tau(C^\bullet,\Gamma)$ denotes the nonzero scalar torsion obtained
from its chirality determinant element under the canonical map
$\det C^\bullet\to\det H^\bullet(C)=\mathbb C$, with the convention
of \cite{CD}. This is the finite-complex torsion, distinct from the
analytic torsion $T_\chi(\Sigma)$. Here
$\det C^\bullet=\bigotimes_k(\det C^k)^{(-1)^k}$, where the
exponent $-1$ denotes the dual line.

Choose a reference time $\tau_*$, write
\(C=C^\bullet_{\theta,\tau_*}\), and define the cochain transport
$S_\tau:C\to C^\bullet_\tau$ and the transported chirality
$\widetilde\Gamma_\tau$ by
\[
 S_\tau=\Pi_{\theta,\tau}|_C,
 \qquad
 \widetilde\Gamma_\tau=S_\tau^{-1}\Gamma_\tau S_\tau.
\]
The maps \(S_\tau\) are locally invertible cochain maps because
\(d_\theta\Pi_{\theta,\tau}=\Pi_{\theta,\tau}d_\theta\).
Thus the transported differential is constant:
\begin{equation}\label{xp:constant-transported-differential}
 S_\tau^{-1}d_\theta S_\tau=d_\theta|_C.
\end{equation}
At the reference time,
\[
 \Pi\dot\Pi\Pi=0,
 \qquad [\Pi,\Gamma]=0,
 \qquad
 \dot{\widetilde\Gamma}=\Pi\dot\Gamma|_C.
\]
For this fixed acyclic complex the chirality determinant element
therefore gives
\begin{equation}\label{xp:chirality-derivative}
 \partial_\tau\log\tau(C^\bullet_\tau,\Gamma_\tau)
   =\tfrac12\operatorname{str}_C
                (\dot{\widetilde\Gamma}\widetilde\Gamma)
   =\sum_{k=0}^{2}(-1)^{k+1}
           \tr_{C^k}(\Pi\Gamma\dot\Gamma).
\end{equation}

To evaluate these traces, define the scalar $h$, the horizontal
two-form $\delta$, and the bundle maps $\mathsf A_j$ on horizontal
degree $j$ for $j=0,1$ by
\begin{equation}\label{xp:chirality-bundle-maps}
 \begin{aligned}
 h=\alpha(\dot X)=-\dot\alpha(X),
 \qquad
 \delta&=d\dot\alpha+(\iota_Xd\dot\alpha)\wedge\alpha,\\
 \mathsf A_0u=2(\ell^2)^{-1}(u\wedge\delta\wedge d\alpha),
 \qquad \mathsf A_1u&=\ell^{-1}(u\wedge\delta).
 \end{aligned}
\end{equation}
Here \(\delta\) is horizontal.  Differentiating
\eqref{xp:contact-chirality} gives the following diagonal blocks
of \(\Gamma\dot\Gamma\):
\[
 \renewcommand{\arraystretch}{1.2}
 \begin{array}{ccc}
 \toprule
 k&U^{k-1}\wedge\alpha&U^k\\ \midrule
 0&0&\mathsf A_0-h\\
 1&\mathsf A_0+h&\mathsf A_1-h\\
 2&\mathsf A_1+h&-h\\ \bottomrule
 \end{array}
\]
The remaining blocks interchange the summands and have zero trace
after compression by \(\Pi\).  Define
\[
 a_j=\tr_{U^j}(\Pi \mathsf A_j)\quad(j=0,1),
 \qquad h_j=\tr_{U^j}(\Pi h)\quad(0\leq j\leq4).
\]
Writing \(\gamma_k=\tr_{C^k}(\Pi\Gamma\dot\Gamma)\), we obtain
\begin{align}\label{xp:chirality-trace-cancellation}
 \gamma_0&=a_0-h_0,\\
 \gamma_1&=a_0+h_0+a_1-h_1,\notag\\
 \gamma_2&=a_1+h_1-h_2,\notag\\
 -\gamma_0+\gamma_1-\gamma_2
      &=2h_0-2h_1+h_2.\notag
\end{align}
On the other hand, the diagonal of
\(\alpha\wedge\iota_{\dot X}\) is multiplication by \(h\)
on \(U^{k-1}\wedge\alpha\) and zero on \(U^k\).  Thus
\begin{equation}\label{xp:full-cut-supertrace}
 H=\operatorname{str}_C(\Pi\alpha\wedge\iota_{\dot X})
   =-h_0+h_1-h_2+h_3-h_4.
\end{equation}
The cut Lefschetz isomorphisms intertwine both \(\Pi\) and
multiplication by \(h\), so \(h_3=h_1\) and \(h_4=h_0\).
Combining \eqref{xp:chirality-derivative}--\eqref{xp:full-cut-supertrace}
proves
\begin{equation}\label{torsionvariation}
 \partial_\tau\log\tau(C^\bullet_\tau,\Gamma_\tau)=-H.
\end{equation}
The computation requires neither a nilpotent compressed operator nor
a cut containing only the zero resonance.

Put $q=(\dim M-1)/2=2$. For the full ranks in \eqref{ranks},
the parity exponent $Q_{\rm cut}$ in the sign of
\cite[(6.6)]{CD} is constant:
\[
 Q_{\rm cut}
    =\sum_{k=0}^{2}(-1)^k(3-k)\dim C^k
    =3-6+6=3.
\]
The finite-cut formula \cite[(6.6)]{CD}, with $q=2$, reads
\begin{equation}\label{dynamicalbridge}
 \tau_{\rm dyn}(\theta,\tau)
 =(-1)^{Q_{\rm cut}}G(0,\theta,\tau)
   \tau(C^\bullet_{\theta,\tau},\Gamma_\tau).
\end{equation}
For the fixed acyclic character, contact deformation invariance and
\eqref{torsionvariation} therefore give
\[
 0=\partial_\tau\log\tau_{\rm dyn}
   =\partial_\tau\log G(0,\theta,\tau)-H(\theta,\tau),
\]
and hence
\[
 \frac{G(0,\theta,\tau)}{G(0,\theta,0)}
        =\exp\!\left(\int_0^\tau H(\theta,v)\,dv\right).
\]

Set \(\mathcal K_\tau=\alpha_\tau\wedge\iota_{\dot X_\tau}\).
Using the fixed sources $f_j$ and dual tests $\ell_i$ of
\eqref{xp:frame-matrices}, define in each full degree
\[
 F_{ij}(z,\tau)=\ell_i(\Pi_{z,\tau}f_j),
 \qquad
 J_{ij}(z,\tau)=\ell_i
       (\Pi_{z,\tau}\mathcal K_\tau\Pi_{z,\tau}f_j).
\]
Write $F_k,J_k$ for these matrices in degree $k$. The coordinate
matrix of the compressed bundle map is \(F^{-1}J\). Its continuity
follows from the sequence of mappings
\[
 \mathcal H_2
 \xrightarrow{\ \Pi_{z,\tau}\ }\mathcal H_1
 \xrightarrow{\ \mathcal K_\tau\ }\mathcal H_1
 \xrightarrow{\ \Pi_{z,\tau}\ }\mathcal H_0
 \xrightarrow{\ \ell_i\ }\mathbb C.
\]
The matrices \(F\) remain uniformly invertible. Define the
extension of the finite trace to complex $z$ by
\[
 H(z,\tau)=\sum_{k=0}^{5}(-1)^k
              \tr\bigl(F_k(z,\tau)^{-1}J_k(z,\tau)\bigr)
\]
This extends continuously across \(z=0\). Choose a uniform bound
$K_H>0$ with $|H|\leq K_H$. Then
\[
 \left|\frac{G(0,\theta,\tau)}{G(0,\theta,0)}-1\right|
    \leq \exp\!\left(\int_0^{|\tau|}|H(\theta,v)|\,dv\right)-1
    \leq e^{K_H|\tau|}-1.
\]
Only the finite trace is extended to the nonacyclic character;
the scalar acyclic torsion identity has been used at \(\theta\neq0\).
\end{proof}

\subsection{Specialization at the crossing}\label{sec:specialization}
\begin{theorem}\label{obstruction}
For sufficiently small real $\theta\neq0$, choose any crossing
$\sigma_\theta$ furnished by \eqref{cross}. At that metric, either
$Z(s,\theta,\sigma_\theta)$ has a simple zero at $s=0$, or it is
regular and nonzero there and
\begin{equation}\label{smallvalue}
 \frac{Z(0,\theta,\sigma_\theta)}{Z(0,\theta,0)}\longrightarrow0.
\end{equation}
In either case the ordinary assertion of Conjecture 9 fails at every
sufficiently small such crossing. The geometric inputs are detailed in the appendix, with
explicit attribution to the primary spectral formulas used there.
\end{theorem}
\begin{proof}
At the hyperbolic endpoint \eqref{hyper}, \eqref{qhyper}, and
\eqref{divisor} give the exact value
\begin{equation}\label{hypervalue}
 C_\theta:=Z(0,\theta,0)
       =-G(0,\theta,0)c_0(\theta)^2\neq0.
\end{equation}
At a crossing, \eqref{two} gives $e_2=0$, so $Q(s)=s(s-e_1)$.
In the formulas at this crossing, $c,e_1,e_2$, and $Q$ are evaluated
at $(\theta,\sigma_\theta)$ when their parameter arguments are omitted.
Fix $(\theta,\sigma_\theta)$ first and cancel the common $s$ factor
in \eqref{divisor}. In both multiplicity cases this gives
\begin{equation}\label{specialized}
 \begin{aligned}
 Z(s,\theta,\sigma_\theta)
   &=G(s,\theta,\sigma_\theta)(s-c)(s-e_1),\\
 Z(0,\theta,\sigma_\theta)
   &=G(0,\theta,\sigma_\theta)c e_1.
 \end{aligned}
\end{equation}
If $m=3$, then $e_1=0$ and the first line is $G(s)(s-c)s$. Its
zero at the origin is simple because $c\ne0$ and $G(0)\ne0$.
If $m=2$, then $e_1\ne0$ and the value in the second line is nonzero.

Divide the second line of \eqref{specialized} by \eqref{hypervalue}:
\begin{equation}\label{exactratio}
 \frac{Z(0,\theta,\sigma_\theta)}{C_\theta}
 =-\frac{G(0,\theta,\sigma_\theta)}{G(0,\theta,0)}
   \frac{c(\theta,\sigma_\theta)}{c_0(\theta)}
   \frac{e_1(\theta,\sigma_\theta)}{c_0(\theta)}.
\end{equation}
This identity includes the simple-zero case. The first factor tends
to one by Lemma~\ref{cutvariation} and $\sigma_\theta=O(\theta^2)$.
The other two limits follow directly from the quotient estimates:
\begin{equation}\label{coefficientratios}
 \frac{c(\theta,\sigma_\theta)}{c_0(\theta)}
   =\frac{C(\theta,\sigma_\theta)}{C(\theta,0)}\longrightarrow1,
 \qquad
 \frac{e_1(\theta,\sigma_\theta)}{c_0(\theta)}
   =\frac{A(\theta,\sigma_\theta)}{C(\theta,0)}\longrightarrow0.
\end{equation}
Here $C(\theta,0)\to-\kappa/2\ne0$ and
$A(\theta,\sigma_\theta)\to A(0,0)=0$. This proves
\eqref{smallvalue} and the three estimates in \eqref{threeestimates}.

Fried's hyperbolic theorem~\cite{Fried} identifies
$|C_\theta|=T_{\chi_\theta}(\Sigma,g_0)^2$. Acyclic unitary analytic
torsion is independent of the metric on this odd-dimensional base,
so this is also the target at $g_{\sigma_\theta}$. A simple zero
contradicts the required nonvanishing; in the second case the modulus
ratio is less than one for all sufficiently small $\theta$, contradicting
the required value. The full zero complex, rather than horizontal
exactness or semisimplicity, supplied the exhaustive two cases.
\end{proof}

\begin{proof}[Proof of Theorem~\ref{mainresult}]
Use the base constructed in the appendix and any sequence of
nonzero real $\theta_j\to0$. Choose the crossings
$\sigma_{\theta_j}=O(\theta_j^2)$ from \eqref{cross}, and put
$g_j=g_{\sigma_{\theta_j}}$, $\chi_j=\chi_{\theta_j}$.
In both cases of \eqref{two}, $Q(s)=s(s-e_1)$ at the crossing,
so \eqref{divisor} is the holomorphic product
$G(s)(s-c)(s-e_1)$. Its value is $G(0)c e_1$, including the case
$e_1=0$. Equation \eqref{exactratio} therefore holds in both cases.
The three factor limits in Theorem~\ref{obstruction}, and the
hyperbolic torsion identity, give the assertion. Smooth convergence
of the metrics follows from $\sigma_{\theta_j}\to0$.
\end{proof}

\begin{remark}[Why a fixed complex disk is essential]\label{sec:model}
The quantitative ingredient is uniform holomorphy on a fixed complex
character disk. Real $C^0$ continuity, evenness, the two vanishing
axes, a real crossing, and divisor cancellation do not give the relative
estimate. A concrete coefficient model on $\tau\geq0$ is
\[
 a=\theta^2,\quad c=-a,\quad t=\tau-a,\quad
 e_1=\frac{2a\tau}{\tau+a},\quad e_2=a(\tau-a),\quad G=1,
\]
with $e_1(0,0)=0$. It has the required axis identities, is continuous
and even in real $\theta$, and has its crossing at $\tau=a$.
For $a>0$ and $t\ne0$, the divisor gives
\[
 Z(0,\theta,\tau)
   =\frac{c e_2}{t}=-a^2.
\]
At $\tau=a$ one has $e_1=a$ and $e_2=0$, so the specialized product is
\[
 Z(s,\theta,a)=(s+a)(s-a),\qquad Z(0,\theta,a)=-a^2.
\]
Thus the crossing
alone forces no ordinary value discrepancy. The coefficient is
holomorphic in $\theta$ at each fixed $\tau>0$, but its poles
$\theta=\pm i\sqrt\tau$ approach the origin. It fails exactly the
uniform-disk hypothesis furnished by \eqref{Neumann} for the actual
geometric family.

The present proof retains all geometric spectral inputs but replaces
mixed metric-character differentiability by uniform complex-character
holomorphy and a finite-cut normalization ratio. It neither exchanges
evaluation at zero with meromorphic parameter specialization nor assumes
that ordinary zeta values are continuous through a crossing.
\end{remark}

\appendix
\section{Geometric and spectral inputs}\label{app:inputs}
\subsection{The hyperbolic base and acyclic characters}\label{app:base}

We first construct the base and the acyclic character interval,
then give the conformal and spectral inputs in their source conventions.



Let $S$ be a closed oriented surface of genus two. In a symplectic basis
of its first integral homology $H_1(S;\mathbb Z)$, define the matrices
$A,T$, and the intersection-form matrix $J_{\mathrm{sp}}$ as follows.
Here $\Id_m$ denotes the $m\times m$ identity matrix, and a superscript
$T$ denotes matrix transpose:
\begin{equation}\label{app:matrix}
 A=\begin{pmatrix}2&1\\1&-1\end{pmatrix},\qquad
 T=\begin{pmatrix}A&-\Id_2\\\Id_2&0\end{pmatrix},\qquad
 J_{\mathrm{sp}}=\begin{pmatrix}0&\Id_2\\-\Id_2&0\end{pmatrix}.
\end{equation}
Since $A^T=A$, a block multiplication gives
\begin{equation}\label{app:symplectic}
 T^TJ_{\mathrm{sp}}T
 =\begin{pmatrix}A&\Id_2\\-\Id_2&0\end{pmatrix}
  \begin{pmatrix}\Id_2&0\\-A&\Id_2\end{pmatrix}
 =J_{\mathrm{sp}}.
\end{equation}
Thus $T\in\operatorname{Sp}(4,\mathbb Z)$, the group of integral matrices
preserving $J_{\mathrm{sp}}$. Its characteristic polynomial $p$ is
\begin{align}\label{app:polynomial}
 p(x)&=\det(x\Id_4-T)
      =\det(x^2\Id_2-xA+\Id_2)\notag\\
     &=(x^2-2x+1)(x^2+x+1)-x^2
      =x^4-x^3-x^2-x+1.
\end{align}
Let $\overline p$ denote reduction of $p$ modulo two, let $\Phi_5$ denote
the fifth cyclotomic polynomial, and let $\operatorname{ord}_G(a)$ denote
the multiplicative order of $a$ in a group $G$. Then
\begin{equation}\label{app:irreducible}
 \overline p(x)=x^4+x^3+x^2+x+1=\Phi_5(x),\qquad
 \operatorname{ord}_{(\mathbb Z/5\mathbb Z)^\times}(2)=4.
\end{equation}
Consequently $\overline p$ is irreducible over $\mathbb F_2$, and $p$
is irreducible over $\mathbb Q$. Moreover,
\begin{equation}\label{app:noncyclotomic}
 p(1)=-1,\qquad p(2)=3.
\end{equation}
The resulting real root in $(1,2)$ excludes cyclotomicity. The nonzero
coefficient of $x$ excludes a representation $p(x)=q(x^k)$ for any
polynomial $q$ and integer $k>1$.

The mapping-class action onto $\operatorname{Sp}(4,\mathbb Z)$ is
surjective. The homological criterion of Casson--Bleiler, in the form
of \cite[Propositions 1--2]{MS}, therefore supplies an
orientation-preserving pseudo-Anosov $\varphi:S\to S$ inducing $T$.
Its mapping torus
\begin{equation}\label{app:mappingtorus}
 \Sigma=(S\times[0,1])/((x,1)\sim(\varphi(x),0))
\end{equation}
is a closed oriented hyperbolic three-manifold by fibered
hyperbolization \cite{Thurston}; choose a hyperbolic metric $g_0$ on it.
The homology Wang sequence, using
$\varphi_*|_{H_0(S;\mathbb R)}=\Id$, gives
\begin{equation}\label{app:wang}
 0\longrightarrow\operatorname{coker}(T-\Id_4)
 \longrightarrow H_1(\Sigma;\mathbb R)
 \longrightarrow\mathbb R\longrightarrow0.
\end{equation}
Here the cokernel is taken on $H_1(S;\mathbb R)$. By
\eqref{app:noncyclotomic}, $\det(\Id_4-T)=p(1)\ne0$, so the
cokernel is zero. Since an endomorphism of a finite-dimensional space
has equal kernel and cokernel dimensions, we obtain
\begin{equation}\label{app:b1}
 b_1(\Sigma)=1+\dim\ker(T-\Id_4)=1.
\end{equation}



Fix a nonzero real harmonic one-form $\omega$ on $(\Sigma,g_0)$. For real
$\theta$, let $d_\theta^*$ be the $L^2(g_0)$ adjoint of the twisted
differential $d_\theta$, and define the degree-$k$ Hodge Laplacian by
\begin{equation}\label{app:hodgeoperators}
 d_\theta=d+i\theta\omega\wedge,
 \qquad
 \Delta_{\mathrm H,\theta}^{(k)}
 =d_\theta^*d_\theta+d_\theta d_\theta^*
 \quad\text{on $k$-forms}.
\end{equation}
The connection is unitary when $\theta\in\mathbb R$. Choose a loop
$\gamma$ with $I_\gamma=\int_\gamma\omega\ne0$. If
\begin{equation}\label{app:nontrivialcharacter}
 0<|\theta|<\frac{2\pi}{|I_\gamma|},
\end{equation}
then $e^{i\theta I_\gamma}\ne1$. Changing the convention for parallel
transport replaces $\theta$ by $-\theta$ and does not affect this
nontriviality argument.

Write $b_k=\dim H^k(\Sigma;\mathbb R)$ for the degree-$k$ Betti number.
At $\theta=0$, the kernels of the scalar and one-form Hodge
Laplacians have dimensions $b_0=b_1=1$. Choose $\eta>0$ below their
first positive eigenvalues, and write $\mathbf1_{[0,\eta)}(D)$ for the
spectral projection of a self-adjoint operator $D$ onto that interval.
Smooth elliptic perturbation theory, after
shrinking $\eta$ and the character interval if necessary, gives
\begin{equation}\label{app:hodgeprojection}
 \operatorname{rank}\mathbf1_{[0,\eta)}
       (\Delta_{\mathrm H,\theta}^{(0)})
 =\operatorname{rank}\mathbf1_{[0,\eta)}
       (\Delta_{\mathrm H,\theta}^{(1)})=1.
\end{equation}
Write $\mu_0(\theta)$ for the eigenvalue in the scalar slot, and
choose an eigenfunction $f_\theta$ with $\|f_\theta\|_2=1$, where
$\|\cdot\|_2$ is the $L^2(g_0)$ norm on the base.
A scalar harmonic section is parallel. Nontrivial holonomy therefore
implies
\begin{equation}\label{app:positiveeigenvalue}
 \ker\Delta_{\mathrm H,\theta}^{(0)}=0,
 \qquad 0<\mu_0(\theta)<\eta
 \quad(0<|\theta|\ll1).
\end{equation}
The differential commutes with the Hodge Laplacians, so
\begin{align}\label{app:exactslot}
 \|d_\theta f_\theta\|_2^2&=\mu_0(\theta)>0,\notag\\
 \Delta_{\mathrm H,\theta}^{(1)}d_\theta f_\theta
 &=d_\theta\Delta_{\mathrm H,\theta}^{(0)}f_\theta
  =\mu_0(\theta)d_\theta f_\theta.
\end{align}
This nonzero exact form fills the unique small one-form slot in
\eqref{app:hodgeprojection}. In particular,
\begin{equation}\label{app:h01}
 H^0(\Sigma;\chi_\theta)=H^1(\Sigma;\chi_\theta)=0.
\end{equation}
Unitary Hodge duality, or Poincar\'e duality applied also to the
inverse character $\chi_{-\theta}$, gives
\begin{equation}\label{app:acyclic}
 H^k(\Sigma;\chi_\theta)=0\quad(0\leq k\leq3),
 \qquad 0<|\theta|\ll1.
\end{equation}

The same argument gives the gap needed for the weight-one Selberg
factor. When $\theta\ne0$ is small, acyclicity identifies coclosed
one-forms with the coexact sector. Every one-form eigenvalue below
$\eta$ is exact. Thus
\begin{equation}\label{app:coexactgap}
 \operatorname{Spec}\!\left(
 \Delta_{\mathrm H,\theta}^{(1)}\big|_{\ker d_\theta^*}\right)
 \subset[\eta,\infty),\qquad 0<|\theta|\ll1.
\end{equation}
This is a gap on one-forms in the indicated sector; the scalar
eigenvalue in \eqref{app:positiveeigenvalue} tends to zero.

For completeness, the quadratic coefficient of that scalar eigenvalue
can be read directly from the scalar operator. Let
$V=\operatorname{vol}_{g_0}(\Sigma)$ and $f_0=V^{-1/2}$. Harmonicity
gives $d^*\omega=0$. Write $K_1,K_2$ for the linear and quadratic
coefficients in the scalar-operator expansion:
\begin{equation}\label{app:scalarexpansion}
 \Delta_{\mathrm H,\theta}^{(0)}
 =\Delta_{\mathrm H,0}^{(0)}+\theta K_1+\theta^2K_2,
 \qquad
 K_1f=-2i\langle\omega,df\rangle_{g_0},\qquad
 K_2f=|\omega|^2f.
\end{equation}
In particular $K_1f_0=0$. Primes in the next two equations denote
$\theta$-derivatives at $\theta=0$. Differentiating the simple eigenvalue
equation at zero gives
\begin{equation}\label{app:scalarfirstderivative}
 \mu_0'(0)=\langle f_0,K_1f_0\rangle=0,
 \qquad
 \Delta_{\mathrm H,0}^{(0)}f_0'=0.
\end{equation}
Choose the phase of $f_\theta$ so that $\langle f_0',f_0\rangle=0$.
Then $f_0'=0$, and the second derivative gives
\begin{equation}\label{app:scalarsecondderivative}
 \frac12\mu_0''(0)=\langle f_0,K_2f_0\rangle
 =\frac{\|\omega\|_2^2}{V}=:\kappa>0.
\end{equation}
Complex conjugation identifies the scalar operators at $\theta$ and
$-\theta$. The simple eigenvalue is therefore even, and analytic
elliptic perturbation theory yields
\begin{equation}\label{app:scalarquadratic}
 \mu_0(\theta)=\kappa\theta^2+O(\theta^4),\qquad
 -1+\sqrt{1-\mu_0(\theta)}
 =-\frac{\kappa}{2}\theta^2+O(\theta^4).
\end{equation}

Finally, for the unit sphere bundle $\pi:M\to\Sigma$, the fiber is
$S^2$. Let $E_2^{p,q}$ be the second page of its local-coefficient
Leray--Serre spectral sequence, with $p$ the base degree and $q$ the
fiber degree. It has
\begin{equation}\label{app:spherepullback}
 E_2^{p,q}
 =H^p(\Sigma;\chi_\theta\otimes H^q(S^2;\mathbb C))
 =0,
 \qquad q\in\{0,2\},
\end{equation}
and no other nonzero rows. The orientation of the bundle makes the
fiber cohomology local systems trivial. Hence
$H^\bullet(M;\pi^*\chi_\theta)=0$.

\subsection{The nondegenerate conformal direction}\label{app:conformal}

In this subsection $\Delta_{g_0}$ denotes the nonpositive scalar
Laplacian used by CDDP, and $X=X_0$, $\alpha=\alpha_0$. Write tangent
vectors as $\xi=(\xi_H,\xi_V)$ and $\eta=(\eta_H,\eta_V)$ in the
Levi--Civita horizontal/vertical splitting. Define the tangent-bundle
rotation $\mathsf I$ and invariant two-form $\psi$ by
\[
 \mathsf I_{(x,v)}(\xi_H,\xi_V)
   =(v\times\xi_V,v\times\xi_H),\qquad
 \psi(\xi,\eta)=d\alpha_0(\mathsf I\xi,\eta).
\]
Here $\times$ is the cross product determined by $g_0$ and the base
orientation; $\psi$ is CDDP's additional invariant two-form.
Let $J(x,v)=(x,-v)$ be velocity reversal. Here ``outgoing'' denotes
the resonant cut selected by the anisotropic spaces fixed above,
and ``incoming'' denotes its image under $J^*$. Define the closed
horizontal zero one-forms by
$\mathcal C_{\mathrm{cl}}=\cU^1_{0,0}\cap\ker d$, and set
$\mathcal C_\psi=\{v\circ\mathsf I:v\in\mathcal C_{\mathrm{cl}}\}$,
where composition means $(v\circ\mathsf I)(\xi)=v(\mathsf I\xi)$.
Let $\pi_*$ denote fiber
integration with CDDP's fiber-orientation convention.
Normalize a real nonclosed $u\in\mathcal C_\psi$ by
\begin{equation}\label{app:unormalization}
 \pi_*(d\alpha_0\wedge u)=\omega,
 \qquad u_*=J^*u,
 \qquad N=\|\omega\|_2^2.
\end{equation}
Their Lemma 3.17 and Corollary 5.1 identify the corresponding
incoming pushforward with the same $\omega$. Define the distribution
$F$ and the smooth function $f$ by
\begin{equation}\label{app:Ff}
 \pi_*(\alpha_0\wedge du\wedge du_*)
   =F\,d\operatorname{vol}_{g_0},
 \qquad f=\Delta_{g_0}|\omega|^2.
\end{equation}
For a smooth function $h$ on $\Sigma$, let $\widetilde h$ be its
invariant lift to the universal cover $\mathbb H^3$, and let
$d_{\mathbb H^3}$ denote hyperbolic distance. CDDP's smoothing
operator $Q_4$ is defined by the hyperbolic-distance kernel
\begin{equation}\label{app:Q4definition}
 \widetilde{Q_4h}(x)
  =\int_{\mathbb H^3}
    \frac{\widetilde h(y)}{\cosh^4 d_{\mathbb H^3}(x,y)}
      \,d\operatorname{vol}_{\mathbb H^3}(y).
\end{equation}
It descends to $\Sigma$ and extends to distributions
\cite[\S5.1.2]{CDDP}. Their Theorem 5 and Proposition A.1 give
\begin{equation}\label{app:Q4identity}
 Q_4F=-\frac16f,\qquad f\ne0.
\end{equation}
Indeed, $f=0$ would make $|\omega|^2$ constant on the compact base,
which is excluded for a nonzero harmonic one-form by the cited
proposition. The operator $Q_4$ has a symmetric periodized smoothing
kernel, so it is self-adjoint in the distribution--test-function pairing.
Set
\begin{equation}\label{app:bchoice}
 b=Q_4f.
\end{equation}
With $\langle F,b\rangle$ denoting the action of the distribution $F$
on the test function $b$ relative to $d\operatorname{vol}_{g_0}$, we have
\begin{equation}\label{app:basepairing}
 \langle F,b\rangle
 =\langle Q_4F,f\rangle
 =-\frac16\|f\|_2^2\ne0.
\end{equation}

The sign in the passage back to $M$ is essential. The orientation
induced by CDDP's fiber integration is opposite to the contact
orientation $\alpha_0\wedge(d\alpha_0)^2$. With the latter used for
integrals over $M$, the convention is
\begin{equation}\label{app:orientation}
 \int_M\pi^*a\,\Omega
 =-\int_\Sigma a\,\pi_*\Omega
 \quad\text{for top forms $\Omega$ and smooth $a$},
 \qquad
 \pi_*\big(\alpha_0\wedge(d\alpha_0)^2\big)
 =-8\pi\,d\operatorname{vol}_{g_0}.
\end{equation}
Equivalently, CDDP's equation (5.6), tested against $f$, yields
\begin{equation}\label{app:Spositive}
 S:=\int_M\pi^*b\,\alpha_0\wedge du\wedge du_*
   =-\langle F,b\rangle
   =\frac16\|f\|_2^2>0.
\end{equation}

The normalization of the outgoing--incoming pairing is equally
explicit. For outgoing and incoming horizontal one-forms $v,v_*$, put
\begin{equation}\label{app:Bpairing}
 \mathcal B(v,v_*)
 =\int_M\alpha_0\wedge d\alpha_0\wedge v\wedge v_*.
\end{equation}
Let $\mathcal Iu=u\circ\mathsf I$ be the induced rotation of one-forms,
and let brackets denote de Rham cohomology classes. The identities in
\cite[Lemma~3.5, (3.57), (3.67), and the calculation preceding (3.59)]{CDDP},
together with the normalization
\eqref{app:unormalization}, give
\begin{align}\label{app:Bnormalization}
 u_{\rm cl}:=\mathcal Iu,\qquad
 [u_{\rm cl}]&=-\frac{1}{4\pi}[\pi^*\omega],\qquad
 \pi_*(\psi\wedge u_{\rm cl})=\omega,\notag\\
 \mathcal B(u_{\rm cl},J^*u_{\rm cl})&=-\frac{N}{4\pi},\qquad
 \mathcal B(u,J^*u)=\frac{N}{4\pi}.
\end{align}
The final sign reversal is (3.57); $J$ commutes with $\mathsf I$, so the
same rotation is applied to both entries.

The real horizontal zero one-space is the direct sum of one-dimensional
closed and nonclosed subspaces. CDDP (3.55)--(3.57) make their pairing
matrix diagonal with nonzero diagonal entries. For
$g_\tau=e^{-2\tau b}g_0$, set
$\beta=\partial_\tau\alpha_\tau|_{\tau=0}$. Primes below denote
$\tau$-derivatives at zero on the identified degree-one spectral cut.
The subscript $\mathrm{cut}$ denotes compression to that cut.
Equations (4.22) and (4.38), with
$\iota_X\beta=-b\circ\pi$, give
\begin{equation}\label{app:Pvariationpairing}
 \mathcal B(P'_{\mathrm{cut}}u,u_*)=iS,
 \qquad P=-i\mathcal L_X,
\end{equation}
and zero variation on the closed row and column. Consequently, in
the closed/nonclosed basis,
\begin{equation}\label{app:positivevariation}
 P'_{\mathrm{cut}}=
 \begin{pmatrix}0&0\\0&ir\end{pmatrix},
 \qquad
 B'_{\mathrm{cut}}=-iP'_{\mathrm{cut}}=
 \begin{pmatrix}0&0\\0&r\end{pmatrix},
 \qquad
 r=\frac{S}{N/(4\pi)}
  =\frac{2\pi}{3N}\|f\|_2^2>0.
\end{equation}
This proves the required sign of the companion-root derivative using
the cited compressed-variation formula. When $b_1=1$, the nonzero
entry is the nondegeneracy condition of CDDP's Theorem 3. That theorem
and Corollary 4.1 provide the untwisted zero ranks used in the main
argument. Small smooth conformal perturbations remain negatively curved.

\subsection{Hyperbolic character clusters}\label{app:hyperbolic}

Here $SO(m)$ is the special orthogonal group in dimension $m$.
Let $\mathbf1$ be its trivial representation and let $\nu_l$ and
$\sigma_p$ be the representations of $SO(2)$ on
$\Lambda^l\mathbb C^2$ and on trace-free symmetric tensors of order
$p$, respectively. Thus $\sigma_0=\mathbf1$, $\sigma_1=\nu_1$,
$\nu_1$ is the sum of weights $+1,-1$, and $\sigma_2$ is the sum of
weights $+2,-2$.
Write $Z^D_\nu(a)$ for DGRS's Selberg function with representation
$\nu$ and complex argument $a$, suppressing the fixed twist
$\chi_\theta$. Write $Z^{\mathrm{BO}}_\nu(t)$ for the corresponding
Bunke--Olbrich function in the centered argument $t$. The scalar
center in dimension three is $\rho=1$, so $t=a-1$ and comparison of
the Euler products gives
\begin{equation}\label{app:centers}
 Z^D_\nu(a)=Z^{\mathrm{BO}}_\nu(a-1).
\end{equation}
The representations relevant to the first three horizontal degrees satisfy
\begin{equation}\label{app:weights}
 \nu_0=\nu_2=\mathbf1,
 \qquad
 \nu_1\otimes\nu_1=\mathbf1^{\oplus2}\oplus\sigma_2.
\end{equation}

Write $\Lambda^j$ for the exterior-power representation on
$\Lambda^j\mathbb C^3$, and $\operatorname{Sym}^2_0$ for the trace-free
second symmetric power of $\mathbb C^3$. The virtual $SO(3)$ lifts,
with restriction to $SO(2)$ understood, are
\begin{align}\label{app:virtuallifts}
 \mathbf1&\longleftrightarrow\Lambda^0,\notag\\
 \nu_1&\longleftrightarrow\Lambda^1-\Lambda^0,
  &&\Lambda^1|_{SO(2)}=\mathbf1\oplus\nu_1,\notag\\
 \sigma_2&\longleftrightarrow\operatorname{Sym}^2_0-\Lambda^1,
  &&\operatorname{Sym}^2_0|_{SO(2)}
       =\mathbf1\oplus\nu_1\oplus\sigma_2.
\end{align}
The subscript $0$ denotes trace-free symmetric tensors. Using full
symmetric powers in the last line would introduce an extra scalar
summand. The pair of weights in each line occurs once. The symmetric
Bunke--Olbrich formula is a graded determinant for that pair; it carries
no additional square or square root.

Write $A\doteq B$ when $A/B$ extends to a holomorphic nonvanishing
function in the spectral neighborhood under discussion. The elliptic
determinants below are zeta-regularized in Bunke--Olbrich's convention.
Bunke--Olbrich
\cite[Proposition 4.7]{BO}, in these conventions, gives
\begin{align}\label{app:determinants}
 Z^D_{\mathbf1}(a)&\doteq
  \det\!\left(\Delta_{\mathrm H,\theta}^{(0)}-1+(a-1)^2\right),\notag\\
 Z^D_{\nu_1}(1+t)&\doteq
  \frac{\det(\Delta_{\mathrm H,\theta}^{(1)}+t^2)}
       {\det(\Delta_{\mathrm H,\theta}^{(0)}+t^2)}.
\end{align}
For a small nontrivial real character, $d_\theta$ identifies the entire
scalar spectrum with the exact one-form spectrum. Thus the exact
sector cancels the denominator, including the small eigenvalue
$\mu_0(\theta)$. The remaining determinant is on the coexact sector,
and \eqref{app:coexactgap} excludes zeros or poles for
$|t|<\sqrt\eta$. At the trivial character, the harmonic contribution
is instead
\begin{equation}\label{app:harmoniccancellation}
 \frac{t^{2b_1}}{t^{2b_0}}=t^{2(b_1-b_0)}=1.
\end{equation}
The positive exact spectra still cancel, and the nonharmonic coexact
spectrum has a positive gap. This is why the centered weight-one
factor contributes no small root, including at $\theta=0$.

For the weight-two pair, the virtual lift in
\eqref{app:virtuallifts} has Casimir shift
$c_{\rm Cas}(\sigma_2)=1-2^2=-3$; the subscript distinguishes this
constant from the small resonance root $c$.
Let $\nabla_\theta$ be the Levi--Civita covariant derivative on the
indicated tensor bundle tensored with the flat connection $d_\theta$;
$\nabla_\theta^*\nabla_\theta$ is its nonnegative rough Laplacian.
Denote the relevant Lichnerowicz operators by $\Delta_2$ and $\Delta_1$.
In DGRS Appendix B's normalization these and the centered operators are
\begin{align}\label{app:casimirshifts}
 \Delta_2&=\nabla_\theta^*\nabla_\theta-6,
 &\Delta_2-c_{\rm Cas}(\sigma_2)&=\nabla_\theta^*\nabla_\theta-3
 &&\text{on $\operatorname{Sym}^2_0$},\notag\\
 \Delta_1&=\nabla_\theta^*\nabla_\theta-2,
 &\Delta_1-c_{\rm Cas}(\sigma_2)&=\nabla_\theta^*\nabla_\theta+1
 &&\text{on $\Lambda^1$}.
\end{align}
At the DGRS argument $a=s+2$, the centered parameter is $s+1$.
Hence the numerator and denominator operators are respectively
\begin{equation}\label{app:weighttwooperators}
 \nabla_\theta^*\nabla_\theta-2+2s+s^2
       \quad\text{on $\operatorname{Sym}^2_0$},
 \qquad
 \nabla_\theta^*\nabla_\theta+2+2s+s^2
       \quad\text{on $\Lambda^1$}.
\end{equation}
The hyperbolic Bochner inequality used in DGRS's proof of Proposition
7.7 is
\begin{equation}\label{app:bochner}
 \nabla_\theta^*\nabla_\theta\geq3
        \quad\text{on $\operatorname{Sym}^2_0$}.
\end{equation}
A flat unitary twist introduces no additional curvature term. The
rough Laplacian on $\Lambda^1$ is nonnegative. Thus the two operators
in \eqref{app:weighttwooperators} are invertible whenever
\begin{equation}\label{app:weighttwodisk}
 |2s+s^2|<\frac12,
\end{equation}
uniformly for the real characters under consideration. The weight-two
factor is holomorphic and nonvanishing in this fixed neighborhood.


We now enumerate the factors in DGRS (7.6). In horizontal degree
$k\in\{0,1,2\}$, $l\in\{0,\ldots,k\}$ indexes the exterior-power
splitting, $p\geq0$ the trace-free symmetric-power degree, and $j\geq0$
the radial index. The associated Selberg representation is
$\nu_l\otimes\nu_{k-l}\otimes\sigma_p$, and its spectral shift $h$ is
\begin{equation}\label{app:shift}
 h=2(j-l)+p+2+k.
\end{equation}
The complete list with $h\leq2$ is
\begin{equation}\label{app:lowshifttable}
\renewcommand{\arraystretch}{1.2}
\begin{array}{ccccl}
 \toprule
 k&l&p&j&\text{Selberg factor}\\ \midrule
 0&0&0&0&Z^D_{\mathbf1}(s+2)\\ \midrule
 1&1&0&0&Z^D_{\nu_1}(s+1)\\
 1&1&1&0&Z^D_{\nu_1\otimes\nu_1}(s+2)\\ \midrule
 2&2&0&0&Z^D_{\mathbf1}(s)\\
 2&2&0&1&Z^D_{\mathbf1}(s+2)\\
 2&2&1&0&Z^D_{\nu_1}(s+1)\\
 2&2&2&0&Z^D_{\sigma_2}(s+2)\\
 2&1&0&0&Z^D_{\nu_1\otimes\nu_1}(s+2)\\ \bottomrule
\end{array}
\end{equation}
All remaining shifts satisfy $h\geq3$. If $|s|<1/2$, then
$\operatorname{Re}(s+h)>2$, so each remaining factor lies in its
zero-free Euler-product half-plane. The normally convergent tail in
the cited DGRS factorization is therefore a holomorphic nonvanishing
factor. Normal convergence is used here in addition to the individual
zero-free property.

Let $\mathscr D_k$ denote the individual horizontal dynamical
determinant in degree $k$. Combining
\eqref{app:weights}--\eqref{app:lowshifttable} gives
\begin{align}\label{app:individualdivisors}
 \mathscr D_0(s)&\doteq Z^D_{\mathbf1}(s+2),\notag\\
 \mathscr D_1(s)&\doteq Z^D_{\mathbf1}(s+2)^2,\notag\\
 \mathscr D_2(s)&\doteq
      Z^D_{\mathbf1}(s)Z^D_{\mathbf1}(s+2)^3.
\end{align}
The degree-two exponent three consists of one scalar factor from
$(l,p,j)=(2,0,1)$ and two scalar summands from $(1,0,0)$.
The opposite small root comes solely from $(2,0,0)$.

For the scalar eigenvalue $\mu_0(\theta)$, the two shifted scalar
equations in \eqref{app:determinants} are
\begin{align}\label{app:scalarroots}
 a=s+2:&\qquad
  \mu_0+s^2+2s=0,
  &&s=-1\pm\sqrt{1-\mu_0},\notag\\
 a=s:&\qquad
  \mu_0+s^2-2s=0,
  &&s=1\pm\sqrt{1-\mu_0}.
\end{align}
Set
\begin{equation}\label{app:c0}
 c_0(\theta)=-1+\sqrt{1-\mu_0(\theta)}<0
 \qquad(0<|\theta|\ll1).
\end{equation}
The roots near zero are $c_0$ and $-c_0$ respectively; the other
roots are near $-2$ and $2$. They are simple because the scalar
eigenspace has dimension one and $\mu_0(\theta)<1$. All other
scalar eigenvalues remain separated from zero. A fixed sufficiently
small circle therefore contains precisely these roots.

DGRS identifies the orders of the individual dynamical determinants
with algebraic resonance multiplicities. For the operator
$B=-\mathcal L_X$, whose spectral coordinate is $s$, the monic
small-cluster characteristic polynomials are consequently
\begin{equation}\label{app:clusterpolynomials}
 P_0(s)=s-c_0,
 \qquad P_1(s)=(s-c_0)^2,
 \qquad P_2(s)=(s-c_0)^3(s+c_0).
\end{equation}
At $\theta=0$, DGRS Proposition 7.7 gives the same horizontal ranks
$(1,2,4,2,1)$. Contact Lefschetz supplies degrees three and four.
No semisimplicity assertion for the degree-two cluster is needed.



The Bunke--Olbrich preprint has a reciprocal inconsistency between
its Ruelle Definition 5.1 and the Selberg factorization in Proposition
5.2. We use only its Selberg determinant formula and calibrate the
hyperbolic modulus independently by Fried's theorem~\cite{Fried}.

The Euler-product check is elementary. Let $\ell(\gamma)=\ell_{g_0}(\gamma)$
be the length of a primitive orbit $\gamma$, let $m(\gamma)\in SO(2)$
be its normal rotational holonomy, and write $\chi=\chi_\theta$.
For an iterate $j\geq1$, put
\begin{equation}\label{app:BOq}
 q=e^{-j\ell(\gamma)},\qquad m=m(\gamma)^j\in SO(2).
\end{equation}
The alternating numerator of the Selberg factorization is
\begin{equation}\label{app:BOcancellation}
 1-q\operatorname{tr}(m)+q^2=\det(\Id-qm).
\end{equation}
It cancels the Selberg denominator. The remaining logarithm is
\begin{equation}\label{app:BOdirectlog}
 -\sum_{\gamma\ \mathrm{primitive}}\sum_{j\geq1}
 \frac{\operatorname{tr}(\chi(\gamma)^j)e^{-sj\ell(\gamma)}}{j}
 =\log\prod_{\gamma\ \mathrm{primitive}}
       \det(\Id-\chi(\gamma)e^{-s\ell(\gamma)}),
\end{equation}
the direct Euler product used throughout the article. This
normalization issue affects neither \eqref{app:clusterpolynomials}
nor the spectral reductions above.

\begin{thebibliography}{99}
\bibitem{Zw} M. Zworski,
\emph{Mathematical study of scattering resonances},
Bull. Math. Sci. 7 (2017), 1--85,
\href{https://doi.org/10.1007/s13373-017-0099-4}{DOI},
Conjecture 9.

\bibitem{RS} D. B. Ray and I. M. Singer,
\emph{$R$-torsion and the Laplacian on Riemannian manifolds},
Adv. Math. 7 (1971), no. 2, 145--210,
\href{https://doi.org/10.1016/0001-8708(71)90045-4}{doi:10.1016/0001-8708(71)90045-4}.

\bibitem{Cheeger} J. Cheeger,
\emph{Analytic torsion and the heat equation},
Ann. of Math. (2) 109 (1979), no. 2, 259--322,
\href{https://doi.org/10.2307/1971113}{doi:10.2307/1971113}.

\bibitem{Muller} W. M\"uller,
\emph{Analytic torsion and $R$-torsion of Riemannian manifolds},
Adv. Math. 28 (1978), no. 3, 233--305,
\href{https://doi.org/10.1016/0001-8708(78)90116-0}{doi:10.1016/0001-8708(78)90116-0}.

\bibitem{Muller93} W. M\"uller,
\emph{Analytic torsion and $R$-torsion for unimodular representations},
J. Amer. Math. Soc. 6 (1993), no. 3, 721--753,
\href{https://doi.org/10.2307/2152781}{doi:10.2307/2152781}.

\bibitem{Fried} D. Fried,
\emph{Analytic torsion and closed geodesics on hyperbolic manifolds},
Invent. Math. 84 (1986), no. 3, 523--540,
\href{https://doi.org/10.1007/BF01388745}{doi:10.1007/BF01388745}.

\bibitem{ShenSurvey} S. Shen,
\emph{Analytic torsion and dynamical flow: a survey on the Fried conjecture},
in \emph{Arithmetic $L$-Functions and Differential Geometric Methods},
Progr. Math. 338, Birkh\"auser, Cham, 2021, 247--299,
\href{https://doi.org/10.1007/978-3-030-65203-6_9}{doi:10.1007/978-3-030-65203-6\_9}.

\bibitem{Fried87} D. Fried,
\emph{Lefschetz formulas for flows}, in
\emph{The Lefschetz Centennial Conference, Part III (Mexico City, 1984)},
A. Verjovsky (ed.), Contemp. Math. 58, Part III,
Amer. Math. Soc., Providence, RI, 1987, 19--69,
\href{https://doi.org/10.1090/conm/058.3/893856}{doi:10.1090/conm/058.3/893856}.

\bibitem{Fried95} D. Fried,
\emph{Meromorphic zeta functions for analytic flows},
Comm. Math. Phys. 174 (1995), no. 1, 161--190,
\href{https://doi.org/10.1007/BF02099469}{doi:10.1007/BF02099469}.

\bibitem{MoSt} H. Moscovici and R. J. Stanton,
\emph{$R$-torsion and zeta functions for locally symmetric manifolds},
Invent. Math. 105 (1991), no. 1, 185--216,
\href{https://doi.org/10.1007/BF01232263}{doi:10.1007/BF01232263}.

\bibitem{Shen} S. Shen,
\emph{Analytic torsion, dynamical zeta functions, and the Fried conjecture},
Anal. PDE 11 (2018), no. 1, 1--74,
\href{https://doi.org/10.2140/apde.2018.11.1}{doi:10.2140/apde.2018.11.1},
\href{https://arxiv.org/abs/1602.00664}{arXiv:1602.00664}.

\bibitem{BObook} U. Bunke and M. Olbrich,
\emph{Selberg zeta and theta functions: a differential operator approach},
Math. Research 83, Akademie-Verlag, Berlin, 1995,
ISBN 978-3-05-501690-5.

\bibitem{BO} U. Bunke and M. Olbrich,
\emph{Theta and zeta functions for odd-dimensional locally symmetric
spaces of rank one}, preprint, 1994,
\href{https://arxiv.org/abs/dg-ga/9407012}{arXiv:dg-ga/9407012},
\href{https://arxiv.org/html/dg-ga/9407012v1}{version 1},
Proposition 4.7 and its graded-determinant convention.

\bibitem{Rugh} H. H. Rugh,
\emph{Generalized Fredholm determinants and Selberg zeta functions for
Axiom A dynamical systems},
Ergodic Theory Dynam. Systems 16 (1996), no.~4, 805--819,
\href{https://doi.org/10.1017/S0143385700009111}{doi:10.1017/S0143385700009111}.

\bibitem{GLP} P. Giulietti, C. Liverani, M. Pollicott,
\emph{Anosov flows and dynamical zeta functions},
Ann. of Math. (2) 178 (2013), no.~2, 687--773,
\href{https://annals.math.princeton.edu/2013/178-2/p06}{Publisher text},
\href{https://doi.org/10.4007/annals.2013.178.2.6}{doi:10.4007/annals.2013.178.2.6}.

\bibitem{GLPErr} P. Giulietti, M. Pollicott, C. Liverani,
\emph{Anosov flows and dynamical zeta functions (Errata)},
2022, \href{https://arxiv.org/abs/2203.04917}{arXiv:2203.04917}.

\bibitem{FS} F. Faure, J. Sj\"ostrand,
\emph{Upper bound on the density of Ruelle resonances for Anosov flows},
Comm. Math. Phys. 308 (2011), no.~2, 325--364,
\href{https://doi.org/10.1007/s00220-011-1349-z}{doi:10.1007/s00220-011-1349-z},
\href{https://arxiv.org/abs/1003.0513}{Primary preprint}.

\bibitem{DZ16} S. Dyatlov, M. Zworski,
\emph{Dynamical zeta functions for Anosov flows via microlocal analysis},
Ann. Sci. \'Ec. Norm. Sup\'er. (4) 49 (2016), no.~3, 543--577,
\href{https://numdam.org/articles/10.24033/asens.2290/}{Publisher text},
\href{https://doi.org/10.24033/asens.2290}{doi:10.24033/asens.2290}.

\bibitem{Bon} Y. Guedes Bonthonneau,
\emph{Perturbation of Ruelle resonances and Faure--Sj\"ostrand anisotropic space},
Rev. Un. Mat. Argentina 61 (2020), no.~1, 63--72,
\href{https://doi.org/10.33044/revuma.v61n1a03}{doi:10.33044/revuma.v61n1a03},
\href{https://arxiv.org/abs/1806.08125}{Primary preprint}.

\bibitem{DR} N. V. Dang, G. Rivi\`ere,
\emph{Topology of Pollicott--Ruelle resonant states},
Ann. Sc. Norm. Super. Pisa Cl. Sci. (5) 21 (2020), special issue,
827--871,
\href{https://doi.org/10.2422/2036-2145.201804_010}{doi:10.2422/2036-2145.201804\_010},
\href{https://arxiv.org/abs/1703.08037}{Primary preprint}, Theorem~2.1.

\bibitem{DZ17} S. Dyatlov, M. Zworski,
\emph{Ruelle zeta function at zero for surfaces},
Invent. Math. 210 (2017), no.~1, 211--229,
\href{https://math.berkeley.edu/~zworski/zazi.pdf}{Primary author PDF}.

\bibitem{CP} M. Ceki\'c, G. P. Paternain,
\emph{Resonant spaces for volume-preserving Anosov flows},
Pure Appl. Anal. 2 (2020), no.~4, 795--840,
\href{https://doi.org/10.2140/paa.2020.2.795}{doi:10.2140/paa.2020.2.795},
\href{https://arxiv.org/abs/1911.03881}{Primary preprint}.

\bibitem{CP25} M. Ceki\'c and G. P. Paternain,
\emph{Resonant forms at zero for dissipative Anosov flows},
Geom. Topol. \textbf{29} (2025), no.~7, 3635--3716.
\href{https://doi.org/10.2140/gt.2025.29.3635}{doi:10.2140/gt.2025.29.3635};
\href{https://arxiv.org/abs/2211.06255v2}{arXiv:2211.06255v2}.

\bibitem{DGRS} N. V. Dang, C. Guillarmou, G. Rivi\`ere, and S. Shen,
\emph{The Fried conjecture in small dimensions},
Invent. Math. \textbf{220} (2020), no.~2, 525--579.
\href{https://doi.org/10.1007/s00222-019-00935-9}{doi:10.1007/s00222-019-00935-9};
\href{https://arxiv.org/abs/1807.01189v3}{arXiv:1807.01189v3}.

\bibitem{CDDP} M. Ceki\'c, B. Delarue, S. Dyatlov, and G. P. Paternain,
\emph{The Ruelle zeta function at zero for nearly hyperbolic 3-manifolds},
Invent. Math. \textbf{229} (2022), no.~1, 303--394.
\href{https://doi.org/10.1007/s00222-022-01108-x}{doi:10.1007/s00222-022-01108-x};
\href{https://arxiv.org/abs/2009.08558v3}{arXiv:2009.08558v3}.
Equation references follow arXiv version~v3.

\bibitem{CD} Y. Chaubet and N. V. Dang,
\emph{Dynamical torsion for contact Anosov flows},
Anal. PDE \textbf{17} (2024), no.~8, 2619--2681.
\href{https://doi.org/10.2140/apde.2024.17.2619}{doi:10.2140/apde.2024.17.2619};
\href{https://arxiv.org/abs/1911.09931v2}{arXiv:1911.09931v2}.
\href{https://arxiv.org/html/1911.09931v2}{Author version~v2}
(equation numbering used here).

\bibitem{BS} J.-M. Bismut, S. Shen,
\emph{Anosov vector fields and Fried sections},
Comm. Math. Phys. 406 (2025), article~240,
\href{https://doi.org/10.1007/s00220-025-05400-8}{doi:10.1007/s00220-025-05400-8},
\href{https://webusers.imj-prg.fr/~shu.shen/publication/cmp_BismutShen25.pdf}
{Primary final author PDF}.

\bibitem{HT26} T. Humbert and Z. Tao,
\emph{Generic twisted Pollicott--Ruelle resonances and zeta function at zero},
preprint, \href{https://arxiv.org/abs/2602.12166v2}{arXiv:2602.12166v2},
3 March 2026 (first version posted 12 February 2026 under the title
\emph{Twisted Pollicott--Ruelle resonances and zeta function at zero on surfaces}).

\bibitem{Spilioti25} P. Spilioti,
\emph{Twisted Ruelle zeta function on hyperbolic manifolds and
complex-valued analytic torsion},
J. Inst. Math. Jussieu 24 (2025), no.~3, 813--845,
\href{https://doi.org/10.1017/S1474748024000586}{doi:10.1017/S1474748024000586},
\href{https://arxiv.org/abs/2004.13474}{arXiv:2004.13474}.

\bibitem{ShenComplex} S. Shen,
\emph{Complex valued analytic torsion and dynamical zeta function on
locally symmetric spaces},
Int. Math. Res. Not. IMRN (2023), no.~5, 3676--3745,
\href{https://arxiv.org/abs/2009.03427}{arXiv:2009.03427}.

\bibitem{JZ24} M. J\'ez\'equel and J. Zung,
\emph{Zeta functions and the Fried conjecture for smooth pseudo-Anosov flows},
preprint, \href{https://arxiv.org/abs/2409.17014v2}{arXiv:2409.17014v2},
18 August 2026 (first version posted 25 September 2024).

\bibitem{ShenYu} S. Shen and J. Yu,
\emph{Morse--Smale flow, Milnor metric, and dynamical zeta function},
J. \'Ec. polytech. Math. 7 (2020), 585--607,
\href{https://doi.org/10.5802/jep.154}{doi:10.5802/jep.154}.

\bibitem{MS} D. Margalit and S. Spallone,
\emph{A homological recipe for pseudo-Anosovs},
\href{https://margalit.droppages.net/papers/homcrit.pdf}{Primary author PDF},
Propositions 1--2.

\bibitem{Thurston} W. P. Thurston,
\emph{Hyperbolic structures on 3-manifolds, II: surface groups and
3-manifolds which fiber over the circle},
\href{https://arxiv.org/abs/math/9801045}{Primary preprint}.
\end{thebibliography}
\end{document}
